\chapter{An \texorpdfstring{$(\oo,1)$}--category of conically smooth stratified bordisms}

We now construct an interval Segal space whose associated $(\oo,1)$-category models $d$-dimensional unoriented conically smooth stratified bordisms. The strategy is the direct stratified
analogue of the construction of $\bord_d$ in \S3.2: we encode composable strings of bordisms by a single \emph{embedded ``total bordism''} over the core of an interval object $T\in \Int$, and we obtain the relevant classifying spaces by applying the concordance resolution to an appropriate smooth $\oo$-sheaf.

Before proceeding with definitions, let us sketch the path. Fix a finite-dimensional real vector space $V$. We first define a bifunctor
\begin{equation}
\tilbordstr_{d,V}\colon \man^{\op}\times \Int \to \set
\end{equation}
which assigns to $(M,T)$ a moduli-type set of bounded closed subsets
$W\subset M\times \core(T)\times V$ satisfying boundary-fixing, conically smooth stratified regularity, and fiberwise submersivity conditions. For each fixed $T$, the functor $\tilbordstr_{d,V}(-,T)$ will be shown to satisfy descent for open covers of manifolds, hence to define a (discrete) smooth $\oo$-sheaf. Restricting along the smooth cosimplicial manifold $\AA^{\bullet} \longra \man$ and taking geometric realization produces a functor
\begin{equation}
\bordstr_{d,V}\colon \Int \longra \sset, \qquad
T \longmapsto \left|\tilbordstr_{d,V}(\AA^\bullet,T) \right|.
\end{equation}
The main technical work of this section is to prove that $\bordstr_{d,V}$ is an interval Segal space, and that its homotopy $1$-category recovers the ordinary category $\cS\mathsf{Cob}_d$ of \Cref{ch5:: def: strat bord ord}.

While the formalism parallels the unstratified case, there are three geometric inputs specific to the
stratified setting: 
\begin{enumerate}
    \item We work with stratified \emph{topological} $d$-manifolds equipped with conically smooth atlases. 
    \item We require a stratified analogue of the differential-topological facts used in Section 3.2, notably that proper conically smooth submersions (also) behave like fiber bundles. 
    \item We rely on the existence and connectedness of conically smooth embeddings of the stratified manifolds under consideration into sufficiently large Euclidean spaces. 
\end{enumerate}
With these in place, we define $\tilbordstr_{d,V}$. 



\section{The unoriented stratified bordism category}

The principal goal of this thesis is to construct an $(\oo,1)$-category of stratified bordisms, $\bordstr_d$. It is useful, however, to first consider the ordinary 1-category of stratified bordisms. This will provide us with a model against which we will compare the homotopy category of $\bordstr_d$.

\subsection{Stratified \texorpdfstring{$d$}{d}-manifolds}

In this thesis we are primarily interested in studying bordism categories of conically smooth stratified spaces with additional regularity features. These are designed around the following premise: physically relevant spacetimes in topological quantum field theories are \emph{topological manifolds}, at the very least. Positive codimension strata encode topological defects which are, in principle, themselves governed by a lower-dimensional topological QFT, thereby necessitating similar regularity features, and so on.

Given the construction of the site $\Bsc_{\leq d,d} \hookra \Bsc$, containing all basics of dimension $d$ and of depth bounded by $d$, as outlined in Chapter 4, one might have thought that this sieve captures exactly the singularity types relevant for topological QFTs with topological defects. This intuition fails, however, because being of \emph{pure dimension} $d$ is strictly less restrictive than being of \emph{constant} topological dimension $d$.   

We will now define a pair of $(\oo,1)$-categories inductively: 
\begin{enumerate}
    \item Let $\cS_{-1}$ be the full $(\oo,1)$-subcategory of $\Snglr$ on $\emptyset$ and let $\cS_{\leq -1} = \cS_{-1}$.
    \item Let $\Bsc_{0}^{\cS_{-1}}$ be the full $(\oo,1)$-subcategory of $\Bsc$ on the basic
    \begin{equation}
        U = \RR^{i}\times \sC(L),
    \end{equation}
    such that $i + \big(\mathsf{dim}(L)+1\big) = 0$ and $L\in \cS_{\leq -1}$.
\end{enumerate}
The natural inclusion functor  
\begin{equation}
    \Bsc_{0}^{\cS_{-1}} \xhookrightarrow{\,\,\,\,\,\,} \Bsc
\end{equation}
is a right fibration, since it is a sieve. We can thus consider the $(\oo,1)$-category $\Mfd(\Bsc_{0}^{\cS_{-1}})$ of $\Bsc_{0}^{\cS_{-1}}$-manifolds; these are conically smooth stratified spaces which admit an atlas of basics in $\Bsc_{0}^{\cS_{-1}}$: disjoint unions of points, each finite subset of which is labelled by some element of a poset, together with stratum-preserving injections among them.


Assume now that $\cS_{k-1}$, $\cS_{\leq k-1}$, and $\Bsc^{\cS_{k-1}}_{k}$ have been defined for all $0\leq k \leq d$. We then make the following definitions.
\begin{definition}\label{ch5:: def: sphere sieve}
    Let $\cS_{d}$ be the full $(\oo,1)$-subcategory of $\Mfd(\Bsc_{d}^{\cS_{d-1}})$ on the objects $X\ra P$ such that $X$ is homeomorphic to a $d$-dimensional sphere. 

    We define $\cS_{\leq d}$ to be the full $(\oo,1)$-subcategory of $\Snglr$ whose objects are those of $\cS_{\leq d-1}$ together with those of $\cS_{d}$.
\end{definition}
 

\begin{definition}
    Let $\Bsc_{d+1}^{\cS_{d}}$ be the full $(\oo,1)$-subcategory of $\Bsc$ on the objects $U= \RR^{i}\times \sC(L)$ such that $L\in \cS_{\leq d}$ and
    \begin{equation}
        i + \big(\mathsf{dim}(L)+1\big) = d+1.
    \end{equation}
\end{definition}

\Cref{ch5:: def: sphere sieve} depends on the map $\Bsc_{d}^{\cS_{d-1}} \longra \Bsc$ being a right fibration. We record it as the following lemma.


\begin{lemma}\label{ch5:: lem: sieve of spheres}
    Let $d \geq 0$. The fully faithful inclusion  
    $\Bsc_{d}^{\cS_{d-1}} \longra \Bsc$ 
    is a right fibration of $(\oo,1)$-categories.
\end{lemma} 
\begin{proof}
To check this, it suffices to show that $\Bsc_{d}^{\cS_{d-1}} \subset \Bsc$ is a sieve: in other words, given a basic $U = \RR^{i}\times \sC(L)$ and an open embedding  
\begin{equation}
    \RR^{i}\times \sC(L) \hookra \RR^{k} \times \sC(\SS^{d-k-1})
\end{equation}
where $\SS^{d-k-1} \in \cS_{\leq d-k-1}$, we must show that $U \in \Bsc_{d}^{\cS_{d-1}}$. Since there is a (non-canonical) homeomorphism
\begin{equation}
    \RR^{k} \times \sC(\SS^{d-k-1}) \cong \RR^{d},
\end{equation}
the codomain is topologically a $d$-manifold. Because $\RR^{i}\times \sC(L)$ admits an open embedding into it, the domain must also be a topological $d$-manifold. For a basic of the form $\RR^{i}\times \sC(L)$, this forces the underlying topological space of $L$ to be homeomorphic to $S^{d-i-1}$. Equivalently, we have
\begin{equation}
L \in \cS_{d-i-1}, \text{ with } \mathsf{Ob}(\cS_{d-i-1}) \subset \mathsf{Ob}(\cS_{\leq d}),  
\end{equation}
ending the proof.
\end{proof}

\begin{observation}
Let us, for the sake of exposition, explicate what these $(\oo,1)$-categories consist of for $d$ in the range $0\leq d \leq 2$.

\begin{enumerate}
    \item[\un{\emph{$d = 0$}}:] The $(\oo,1)$-category $\Bsc_{0}^{\cS_-1}$ contains the single basic $\{\RR^{0}\}$, so an object in the $(\oo,1)$-category $\Mfd(\Bsc_{0}^{\cS_{-1}})$ consists of a finite collection of points stratified by some poset $P$, and a morphism consists of a stratum-preserving injection. An example is indicated in the figure below, where the colors encode different strata and the map is manifest from the relative positions:
    \medskip
    \begin{equation}
        \tikzfig{figures/Stcob/SCO2}{\begin{tikzpicture}
	\begin{pgfonlayer}{nodelayer}
		\node [style=none] (0) at (-2, 2) {};
		\node [style=none] (2) at (-2, -2) {};
		\node [style=0 strat] (3) at (-2, 2) {};
		\node [style=0 strat] (5) at (-2, -2) {};
		\node [style=0 strat yellow] (11) at (0, 1) {};
		\node [style=0 strat yellow] (12) at (0, -1) {};
		\node [style=0 strat red] (15) at (2, -2) {};
	\end{pgfonlayer}
\end{tikzpicture}
} \,\,\,\,\, \xhookrightarrow{\qquad} \,\,\,\,\,\tikzfig{figures/Stcob/SCO1}{\begin{tikzpicture}
	\begin{pgfonlayer}{nodelayer}
		\node [style=none] (0) at (-2, 2) {};
		\node [style=none] (1) at (-2, 0) {};
		\node [style=none] (2) at (-2, -2) {};
		\node [style=0 strat] (3) at (-2, 2) {};
		\node [style=0 strat] (4) at (-2, 0) {};
		\node [style=0 strat] (5) at (-2, -2) {};
		\node [style=0 strat yellow] (11) at (0, 1) {};
		\node [style=0 strat yellow] (12) at (0, -1) {};
		\node [style=0 strat red] (13) at (2, 2) {};
		\node [style=0 strat red] (14) at (2, 0) {};
		\node [style=0 strat red] (15) at (2, -2) {};
	\end{pgfonlayer}
\end{tikzpicture}
} 
        \medskip
    \end{equation}
    \item[\un{\emph{$d = 1$}}:]  An object in $\cS_0$ is a stratified 0-sphere $\SS^0 = \big(S^{0}\ra P\big)$ and morphisms are stratum-preserving injections; likewise for $\cS_{\leq 0}$, with the addition of $\emptyset$. An object of $\Bsc_{1}^{\cS_{0}}$ is a basic of one of two types:
    \begin{itemize}
        \item $U = \RR \times \sC(\emptyset)$, where $\emptyset\in \cS_{\leq 0}$:
        \medskip
        \begin{equation}
            U \, = \, \tikzfig{figures/Stcob/SCO3}{\begin{tikzpicture}
	\begin{pgfonlayer}{nodelayer}
		\node [style=none] (4) at (-3, 0) {};
		\node [style=none] (5) at (3, 0) {};
	\end{pgfonlayer}
	\begin{pgfonlayer}{edgelayer}
		\draw [style=geometric] (4.center) to (5.center);
	\end{pgfonlayer}
\end{tikzpicture}
}\quad ;
            \medskip
        \end{equation}
        \item $V= \sC(\SS^0)$, where $\SS^0$ is a stratified 0-sphere in $\cS_{\leq 0}$.
        \medskip
        \begin{equation}
            V \, = \, \tikzfig{figures/Stcob/SCO4}{\begin{tikzpicture}
	\begin{pgfonlayer}{nodelayer}
		\node [style=0 strat] (3) at (0, 0) {};
		\node [style=none] (4) at (-3, 0) {};
		\node [style=none] (5) at (3, 0) {};
	\end{pgfonlayer}
	\begin{pgfonlayer}{edgelayer}
		\draw [style=geometric] (4.center) to (5.center);
	\end{pgfonlayer}
\end{tikzpicture}
}\quad ;
            \medskip
        \end{equation}
    \end{itemize}  
    Morphisms are open embeddings, so $\Hom_{\Bsc_{1}^{\cS_0}}(V,U) = \emptyset$ and $\Hom_{\Bsc_{1}^{\cS_0}}(U,V)$ is classified by two copies of $\Emb(\RR,\RR)$ one for each side of the cone point. Having only these two types of basics to construct the objects of $\Mfd(\Bsc_{1}^{\cS_{0}})$ means that all of them will be topological manifolds of dimension 1 with a conically smooth stratification. This category packs a few interesting features: it contains the diffeomorphism group of the circle, for example; it contains a copy of $\bDelta_{\sf int}$ — the full subcategory of inert morphisms\footnote{Consider the full subcategory of $\Mfd(\Bsc_{1}^{\cS_{0}})$ on the non-compact objects stratified by connected components. Every open embedding is then determined, up to a contractible choice, by the poset map which by definition must be inert (or consecutive).}.
    \medskip  
    \begin{equation}
        \tikzfig{figures/Stcob/SCO5}{\begin{tikzpicture}
	\begin{pgfonlayer}{nodelayer}
		\node [style=none] (4) at (-4, 3) {};
		\node [style=none] (5) at (-4, -3) {};
		\node [style=0 strat] (6) at (-4, 2) {};
		\node [style=0 strat] (7) at (-4, -2) {};
		\node [style=0 strat] (8) at (-4, 0) {};
		\node [style=none] (9) at (-2, 0) {$\sqcup$};
		\node [style=none] (10) at (0, 3) {};
		\node [style=none] (11) at (0, -3) {};
		\node [style=0 strat] (12) at (0, 2) {};
		\node [style=0 strat] (13) at (0, -2) {};
	\end{pgfonlayer}
	\begin{pgfonlayer}{edgelayer}
		\draw [style=geometric] (4.center) to (5.center);
		\draw [style=geometric] (10.center) to (11.center);
	\end{pgfonlayer}
\end{tikzpicture}
}\,\,\,\,\, \xhookrightarrow{\qquad} \,\,\,\,\, \tikzfig{figures/Stcob/SCO6}{\begin{tikzpicture}
	\begin{pgfonlayer}{nodelayer}
		\node [style=none] (4) at (-2, 3) {};
		\node [style=none] (5) at (-2, -3) {};
		\node [style=0 strat] (6) at (-2, 1) {};
		\node [style=0 strat] (7) at (-2, -2) {};
		\node [style=0 strat] (8) at (-2, 0) {};
		\node [style=none] (9) at (0, 0) {$\sqcup$};
		\node [style=none] (10) at (5, 0) {};
		\node [style=none] (11) at (2, 0) {};
		\node [style=0 strat] (14) at (-2, 2) {};
		\node [style=0 strat] (15) at (-2, -1) {};
		\node [style=none] (16) at (-2, 2.5) {};
		\node [style=none] (17) at (-2, -0.5) {};
		\node [style=0 strat] (18) at (2.5, -1.2) {};
		\node [style=0 strat] (19) at (2.55, 1.225) {};
		\node [style=0 strat] (20) at (4.5, 1.2) {};
		\node [style=0 strat] (21) at (3.575, -1.525) {};
		\node [style=none] (22) at (3.75, 1.525) {};
		\node [style=none] (23) at (3, -1.475) {};
	\end{pgfonlayer}
	\begin{pgfonlayer}{edgelayer}
		\draw [style=geometric] (4.center) to (5.center);
		\draw [style=geometric] (11.center)
			 to [in=-270, out=90, looseness=1.75] (10.center)
			 to [in=270, out=270, looseness=1.75] cycle;
		\draw [style=defect line 1 dark] (17.center) to (16.center);
		\draw [style=defect line 2 dark, in=-15, out=-15, looseness=1.75] (22.center) to (23.center);
	\end{pgfonlayer}
\end{tikzpicture}
}
        \medskip
    \end{equation}
    On the left (the domain of the open embedding) we observe a pair of copies of $\RR$ stratified by three and two $0$-dimensional strata, respectively, and on the right (the codomain of the embedding) we observe another copy of $\RR$, this time stratified by five $0$-dimensional strata, together with a copy of $\SS^1$. 
    \item[\un{\emph{$d = 2$}}:]  An object of $\cS_1$ is a stratified 1-sphere $\SS^1 = \big(S^1 \ra P\big)$ and morphisms are stratum-preserving injections among them; the $(\oo,1)$-category $\cS_{\leq 1}$ contains $\emptyset$, all stratified $0$-spheres $\SS^0$, and all stratified 1-spheres $\SS^1$. Objects in $\Bsc_{2}^{\cS_1}$ are basics of one of the following types:
    \begin{itemize}
        \item $U = \RR^2 \times \sC(\emptyset)$, where $\emptyset \in \cS_{\leq 1}$:
        \medskip
        \begin{equation} 
            \tikzfig{figures/Stcob/SCO7}{\begin{tikzpicture}
	\begin{pgfonlayer}{nodelayer}
		\node [style=none] (0) at (-2.5, 2.5) {};
		\node [style=none] (1) at (2.5, 2.5) {};
		\node [style=none] (2) at (-2.5, -2.5) {};
		\node [style=none] (3) at (2.5, -2.5) {};
	\end{pgfonlayer}
	\begin{pgfonlayer}{edgelayer}
		\draw [style=defect patch 2] (0.center)
			 to (1.center)
			 to (3.center)
			 to (2.center)
			 to cycle;
		\draw [style=ts-dashed] (0.center) to (2.center);
		\draw [style=ts-dashed] (2.center) to (3.center);
		\draw [style=ts-dashed] (3.center) to (1.center);
		\draw [style=ts-dashed] (1.center) to (0.center);
	\end{pgfonlayer}
\end{tikzpicture}
} \quad ;
            \medskip
        \end{equation}
         \item $V = \RR^1 \times \sC(\SS^0)$, where $\SS^0 \in \cS_{\leq 1}$:
        \medskip
         \begin{equation} 
            \tikzfig{figures/Stcob/SCO8}{\begin{tikzpicture}
	\begin{pgfonlayer}{nodelayer}
		\node [style=none] (0) at (-2.5, 2.5) {};
		\node [style=none] (1) at (2.5, 2.5) {};
		\node [style=none] (2) at (-2.5, -2.5) {};
		\node [style=none] (3) at (2.5, -2.5) {};
		\node [style=none] (4) at (0, 2.5) {};
		\node [style=none] (5) at (0, -2.5) {};
	\end{pgfonlayer}
	\begin{pgfonlayer}{edgelayer}
		\draw [style=defect patch 2] (0.center)
			 to (1.center)
			 to (3.center)
			 to (2.center)
			 to cycle;
		\draw [style=ts-dashed] (0.center) to (2.center);
		\draw [style=ts-dashed] (2.center) to (3.center);
		\draw [style=ts-dashed] (3.center) to (1.center);
		\draw [style=ts-dashed] (1.center) to (0.center);
		\draw [style=geometric] (4.center) to (5.center);
	\end{pgfonlayer}
\end{tikzpicture}
} \quad ;
            \medskip
        \end{equation}
         \item $W = \sC(\SS^1)$, where $\SS^1 \in \cS_{\leq 1}$:
        \medskip
         \begin{equation}    
            \tikzfig{figures/Stcob/SCO9}{\begin{tikzpicture}
	\begin{pgfonlayer}{nodelayer}
		\node [style=none] (0) at (-2.5, 0) {};
		\node [style=none] (1) at (2.5, 0) {};
		\node [style=none] (2) at (0, 0) {};
		\node [style=0 strat] (3) at (0, 0) {};
		\node [style=none] (4) at (-1.975, 1.7) {};
		\node [style=none] (5) at (0, -2.6) {};
		\node [style=none] (6) at (1.8, 1.825) {};
	\end{pgfonlayer}
	\begin{pgfonlayer}{edgelayer}
		\draw [style=defect patch 2] (1.center)
			 to [in=270, out=-90, looseness=1.75] (0.center)
			 to [in=90, out=90, looseness=1.75] cycle;
		\draw [style=ts-dashed, in=90, out=90, looseness=1.75] (0.center) to (1.center);
		\draw [style=ts-dashed, in=270, out=270, looseness=1.75] (1.center) to (0.center);
		\draw [style=geometric] (4.center) to (3);
		\draw [style=geometric] (3) to (6.center);
		\draw [style=geometric] (3) to (5.center);
	\end{pgfonlayer}
\end{tikzpicture}
} \quad ;
            \medskip
        \end{equation}
        We chose $\SS^1$ stratified by three $0$-dimensional strata, but there is an infinite collection of such singularity types, one for each way of stratifying the $1$-sphere. This rapid expansion in the number of singularity types when passing from $\Bsc_{1}^{\cS_0}$ to $\Bsc_{2}^{\cS_1}$ helps explain the increasing combinatorial complexity of working with conically smooth topological manifolds. 
        
        Objects in the $(\oo,1)$-category $\Mfd(\Bsc_{2}^{\cS_1})$ are thus topological 2-dimensional manifolds endowed with a maximal conically smooth atlas; the morphisms are conically smooth open embeddings. An example of such an object in this $(\oo,1)$-category is:
        \medskip
        \begin{equation}
            \tikzfig{figures/Stcob/SCO10}{\begin{tikzpicture}
	\begin{pgfonlayer}{nodelayer}
		\node [style=none] (43) at (-2.25, 1.5) {};
		\node [style=none] (44) at (-0.25, -1.5) {};
		\node [style=none] (45) at (4, -1.5) {};
		\node [style=none] (46) at (4, 1.5) {};
		\node [style=none] (28) at (-5, 1.5) {};
		\node [style=none] (29) at (-5, -1.5) {};
		\node [style=none] (30) at (-0.25, -1.5) {};
		\node [style=none] (31) at (-2.25, 1.5) {};
		\node [style=none] (24) at (-5, 1.5) {};
		\node [style=none] (25) at (-4.35, 0.5) {};
		\node [style=none] (26) at (-3, 0.5) {};
		\node [style=none] (27) at (-2.25, 1.5) {};
		\node [style=none] (0) at (-5, 1.5) {};
		\node [style=none] (1) at (-5, -1.5) {};
		\node [style=none] (2) at (-5, 1.5) {};
		\node [style=none] (3) at (-5, -1.5) {};
		\node [style=none] (4) at (4, 1.5) {};
		\node [style=none] (5) at (4, -1.5) {};
		\node [style=none] (6) at (4, 1.5) {};
		\node [style=none] (7) at (4, -1.5) {};
		\node [style=none] (13) at (-0.25, -1.5) {};
		\node [style=0 strat] (14) at (-3, 0.5) {};
		\node [style=0 strat] (15) at (0, -0.75) {};
		\node [style=0 strat] (16) at (2.25, 0.75) {};
		\node [style=none] (17) at (-4.35, 0.5) {};
		\node [style=none] (18) at (-2.25, 1.5) {};
		\node [style=none] (19) at (-4.35, 0.5) {};
		\node [style=none] (20) at (-3, 0.5) {};
		\node [style=none] (21) at (0, -0.75) {};
		\node [style=none] (22) at (-0.25, -1.5) {};
		\node [style=none] (23) at (-5.025, -1.5) {};
		\node [style=none] (32) at (-3, 0.525) {};
		\node [style=none] (33) at (2.25, 0.775) {};
		\node [style=none] (34) at (0, -0.725) {};
		\node [style=none] (35) at (0, -0.75) {};
		\node [style=none] (36) at (-0.25, -1.5) {};
		\node [style=none] (37) at (-3, 0.5) {};
		\node [style=none] (38) at (4.5, 0.75) {};
		\node [style=none] (39) at (2.15, 0.75) {};
		\node [style=none] (40) at (4, -1.5) {};
		\node [style=none] (41) at (4, 1.5) {};
		\node [style=none] (42) at (-2.25, 1.5) {};
	\end{pgfonlayer}
	\begin{pgfonlayer}{edgelayer}
		\draw [style=defect patch b] (46.center)
			 to [in=-180, out=180, looseness=0.75] (45.center)
			 to (44.center)
			 to [in=0, out=180, looseness=0.75] (43.center)
			 to cycle;
		\draw [style=defect patch 2] (28.center)
			 to [in=180, out=-180, looseness=0.75] (29.center)
			 to (30.center)
			 to [in=0, out=180, looseness=0.75] (31.center)
			 to cycle;
		\draw [style=defect patch 2] (27.center)
			 to (24.center)
			 to [in=90, out=0, looseness=0.75] (25.center)
			 to (26.center)
			 to cycle;
		\draw [style=ts-dashed, in=180, out=-180, looseness=0.75] (0.center) to (1.center);
		\draw [style=ts-dashed, in=0, out=0, looseness=0.75] (2.center) to (3.center);
		\draw [style=ts-dashed, in=180, out=-180, looseness=0.75] (4.center) to (5.center);
		\draw [style=ts-dashed, in=0, out=0, looseness=0.75] (6.center) to (7.center);
		\draw [style=geometric, in=-180, out=0] (2.center) to (6.center);
		\draw [style=geometric] (3.center) to (7.center);
		\draw [style=defect patch 1 dark] (6.center) to (7.center);
		\draw [style={open_int_blue_2}] (17.center) to (14);
		\draw [style=defect line 1] (14) to (15);
		\draw [style=defect line 1] (15) to (16);
		\draw [style=defect line 3, in=150, out=0, looseness=0.75] (14) to (16);
		\draw [style=defect line 3] (14) to (18.center);
		\draw [style=defect line 3 dashed, in=-180, out=0, looseness=0.75] (18.center) to (13.center);
		\draw [style=defect line 3, in=-75, out=0, looseness=1.25] (13.center) to (15);
		\draw [style=defect patch 2] (17.center) to (14);
		\draw [style=defect patch 2] (14) to (15);
		\draw [style=defect patch 2] (15) to (13.center);
		\draw [style=defect patch 2] (13.center) to (3.center);
		\draw [style=defect patch 2] (3.center) to (17.center);
		\draw [style=defect patch 2] (19.center)
			 to (20.center)
			 to (21.center)
			 to [in=0, out=-75, looseness=1.25] (22.center)
			 to (23.center)
			 to [in=-75, out=0, looseness=0.75] cycle;
		\draw [style=defect patch 1] (33.center)
			 to [in=0, out=150, looseness=0.75] (32.center)
			 to (34.center)
			 to cycle;
		\draw [style=defect patch b] (42.center)
			 to (41.center)
			 to [in=0, out=0, looseness=0.75] (40.center)
			 to (36.center)
			 to [in=-60, out=0, looseness=1.25] (35.center)
			 to (39.center)
			 to [in=0, out=120, looseness=0.50] (37.center)
			 to cycle;
	\end{pgfonlayer}
\end{tikzpicture}
}
        \end{equation}   
        \medskip
    \end{itemize} 

 \end{enumerate}

\end{observation}

\begin{remark}
    The two $(\oo,1)$-categories $\Bsc_{d}^{\cS_{d-1}}$ and $\Mfd(\Bsc_{d}^{\cS_{d-1}})$ inherit a symmetric monoidal structure given by disjoint union. 
\end{remark}

\begin{remark}
   The preceding discussion of the singularity types encoded by $\Bsc_{d}^{\cS_{d-1}}$ helps clarify the proof of \Cref{ch5:: lem: sieve of spheres}. Namely, it makes manifest that a sufficiently small neighborhood of every point of a basic $U \in \Bsc_{d}^{\cS_{d-1}}$ again lies in $\Bsc_{d}^{\cS_{d-1}}$.
\end{remark}

\begin{definition}
    The $(\oo,1)$-category of topological stratified $d$-manifolds of depth at most $d$ is defined as
    \begin{equation}
        \SMfd_d \, \overset{\sf def}{=} \Mfd(\Bsc_{d}^{\cS_{d-1}}).
    \end{equation}        
    The objects are topological $d$-manifolds endowed with a conically smooth atlas and the morphisms are open embeddings. The $(\oo,1)$-category admits a symmetric monoidal structure given by disjoint union:
    \begin{align}
        \sqcup \maps \SMfd_d \, \times \,  \SMfd_d &\longra \SMfd_d \\
        \big(M \ra P, N \ra Q \big) &\longmapsto (M \sqcup N) \ra (P \sqcup Q).
    \end{align}
\end{definition}




\begin{observation}
By varying the depth of the stratification, we obtain different classes of topological stratified $d$-manifolds. Recall the functor 
\begin{equation}
    \mathsf{depth} \maps \SMfd_d \longra [d].
\end{equation}
A stratified topological $d$-manifold $M$ has depth at most $k$ if for all $x \in M$, the local depth $\mathsf{depth}_x(M) \leq k$. Hence, if $k \leq l$ and $M$ has depth at most $k$, then it also has depth at most $l$.   
\end{observation}

Since we are interested in studying bordisms of stratified $d$-manifolds, we introduce the following broader class of stratified spaces.

\begin{definition}
    Let $\SMfd_{d}^{\partial}$ be the $(\oo,1)$-category which is the pullback 
    \begin{equation}
        \begin{tikzcd}[ampersand replacement=\&]
            {\SMfd_{d}^{\partial}} \&\& {\Snglr^{\partial}} \\
            \\
            {\SMfd_{d}} \&\& \Snglr \quad .
            \arrow[hook, from=1-1, to=1-3]
            \arrow[from=1-1, to=3-1]
            \arrow["\lrcorner"{anchor=center, pos=0.125}, draw=none, from=1-1, to=3-3]
            \arrow["{\mathsf{int}}", from=1-3, to=3-3]
            \arrow[hook, from=3-1, to=3-3]
        \end{tikzcd}
    \end{equation} 
    The $(\oo,1)$-category $\SMfd_d^{\partial}$ contains as objects those conically smooth stratified spaces which admit the structure of a stratified topological $d$-manifold with boundary. The morphisms are open embeddings which respect the boundary. The symmetric monoidal structure is again given by disjoint union.
\end{definition}

\begin{remark}
    It is important to remind the reader that there is a notion of depth and one of interior depth. They are not always in agreement. A smooth manifold with boundary seen as a conically smooth stratified space with $\lag * \rag$-corners has depth 1 and interior depth 0, for example. This highlights the sense in which we think of the boundary structure as an additional structure to that of the underlying stratified space.
\end{remark}

\subsection{Stratified bordisms}

\begin{definition}
Let $d\geq 0$ be an integer, and let $\Sigma_{\sf in}$ and $\Sigma_{\sf out}$ be closed stratified $(d-1)$-manifolds. A \define{stratified bordism} $M\colon \Sigma_{\sf in} \ra \Sigma_{\sf out}$ \define{of dimension} $d$ is a tuple 
\begin{equation}
    \big(M,\varphi_{\sf in},\varphi_{\sf out}\big)
\end{equation}
consisting of the following data:
\begin{enumerate}
    \item A compact stratified topological $d$-manifold with boundary $M$, together with a corner structure $c \maps M \ra \cP(*)$, i.e., a topologically coCartesian map.
    \item Conically smooth open embeddings between stratified $d$-manifolds with boundary:
    \begin{align}
    \varphi_{\sf in} \maps \RR_{\geq 0} \times \Sigma_{\sf in} & \hookra M \\
    \varphi_{\sf out} \maps \RR_{\leq 0} \times \Sigma_{\sf out} & \hookra M 
    \end{align} 
    \item Writing
    \begin{equation}
        \partial_{\sf in}M \,\, \overset{\sf def}{=} \,\, \varphi_{\sf in}\big(\{0\}\times \Sigma_{\sf in}\big), \qquad \partial_{\sf out}M \,\, \overset{\sf def}{=} \,\, \varphi_{\sf out}\big(\{0\}\times \Sigma_{\sf out}\big),
    \end{equation}
    we require that
    \begin{equation}
        \partial M \,=\, \partial_{\sf in}M \sqcup \partial_{\sf out}M.
    \end{equation}
\end{enumerate}    
\end{definition}

As in the smooth case (see \Cref{ch3:: def: iso bordism}), we can define isomorphism classes of stratified bordisms.

\begin{definition}\label{ch5:: def: iso bord str}
    Let 
    \begin{equation}
    (M,\varphi_{\sf in},\varphi_{\sf out}),\qquad (M',\varphi'_{\sf in},\varphi'_{\sf out})
    \end{equation}
    be $d$-dimensional stratified bordisms from $\Sigma_{\sf in}$ to $\Sigma_{\sf out}$.
    An \define{isomorphism of stratified bordisms} is a conically smooth diffeomorphism of topological stratified $d$-manifolds with boundary
    \begin{equation}
    F\colon M \xrightarrow{\;\cong\;} M'
    \end{equation}
    such that
    \begin{align*}
        F\circ \varphi_{\sf in} &= \varphi'_{\sf in},\\
        F\circ \varphi_{\sf out} &= \varphi'_{\sf out}.
    \end{align*}
    We write 
    \begin{equation}
        (M,\varphi_{\sf in},\varphi_{\sf out})\cong (M',\varphi'_{\sf in},\varphi'_{\sf out})
    \end{equation}
    to denote a pair of isomorphic stratified bordisms.
\end{definition}

This definition becomes especially important when we construct the stratified bordism category. Next, we observe that, just as smooth cobordism forms an equivalence relation, so does stratified cobordism. 

\begin{proposition}
    Let $\Sigma_0$ and $\Sigma_1$ be a pair of closed stratified $(d-1)$-manifolds. We say that $\Sigma_0$ and $\Sigma_1$ are \define{stratified cobordant} if there exists a stratified bordism $M\maps \Sigma_0 \ra \Sigma_1$ of dimension $d$. The relation of being stratified cobordant is an equivalence relation.
\end{proposition}

\begin{proof}
    We check reflexivity, symmetry, and transitivity.
    \begin{enumerate}
        \item \textbf{Reflexivity:} Let $\Sigma$ be a closed stratified $(d-1)$-manifold. Put $M = \Sigma\times[0,1]$ with the product stratification and with the evident boundary (corner) structure $c\maps M\ra \cP(*)$ classifying $\partial M = \Sigma\times\{0,1\}$. Let $\varphi_{\sf in}$ and $\varphi_{\sf out}$ be the evident collars, and define
        \begin{equation}
            \partial_{\sf in}M = \Sigma\times\{0\}, \qquad \partial_{\sf out}M = \Sigma\times\{1\}.
        \end{equation}
        Then $(M,\varphi_{\sf in},\varphi_{\sf out})$ is a stratified bordism $\Sigma\ra \Sigma$.

        \item \textbf{Symmetry:} Let $(M,\varphi_{\sf in},\varphi_{\sf out})$ be a stratified bordism $\Sigma_0\ra \Sigma_1$. Let $r\maps \RR\ra \RR$ be reflection $r(t)=-t$. Define
        \[
            \varphi^{\op}_{\sf in}= \varphi_{\sf out}\circ\big(r|_{\RR_{\ge 0}}\times \id_{\Sigma_1}\big),\qquad
            \varphi^{\op}_{\sf out}= \varphi_{\sf in}\circ\big(r|_{\RR_{\le 0}}\times \id_{\Sigma_0}\big).
        \]
        These are again conically smooth open embeddings of stratified $d$-manifolds with boundary, and by construction they satisfy
        \[
            \varphi^{\op}_{\sf in}\big(\{0\}\times\Sigma_1\big)= \partial_{\sf out}M,
            \qquad
            \varphi^{\op}_{\sf out}\big(\{0\}\times\Sigma_0\big)= \partial_{\sf in}M.
        \]
        Hence, $(M,\varphi^{\op}_{\sf in},\varphi^{\op}_{\sf out})$ is a stratified bordism $\Sigma_1\ra \Sigma_0$.

        \item \textbf{Transitivity:} Let $(M_0,\varphi^0_{\sf in},\varphi^0_{\sf out})\maps \Sigma_0\ra \Sigma_1$ and $(M_1,\varphi^1_{\sf in},\varphi^1_{\sf out})\maps \Sigma_1\ra \Sigma_2$ be stratified bordisms of dimension $d$. Using the collars, choose $\epsilon>0$ and form the glued topological space
        \[
            M_{01}= M_0\cup_{\Sigma_1} M_1,
        \]
        where we glue along the identified copies of $\Sigma_1$ in the outgoing boundary of $M_0$ and the incoming boundary of $M_1$. The collars identify neighborhoods of these boundary components with
        \[
            (-\epsilon,0]\times \Sigma_1 \qquad \text{and} \qquad [0,\epsilon)\times \Sigma_1,
        \]
        so the usual collar-gluing construction endows $M_{01}$ with the structure of a stratified $d$-manifold with boundary. Writing $q_i\maps M_i\ra M_{01}$ for the canonical maps, the incoming and outgoing collars are given by
        \[
            \varphi^{01}_{\sf in}= q_0\circ \varphi^0_{\sf in},
            \qquad
            \varphi^{01}_{\sf out}= q_1\circ \varphi^1_{\sf out}.
        \]

        Thus $(M_{01},\varphi^{01}_{\sf in},\varphi^{01}_{\sf out})$ is a stratified bordism $\Sigma_0\ra \Sigma_2$.
    \end{enumerate}
    Therefore stratified cobordism is an equivalence relation.
\end{proof}

\begin{remark}
    The key input in the transitivity argument is the same one used in the smooth unstratified case: collars turn the common boundary component into a product neighborhood, which makes gluing possible and canonical up to isomorphism. The only additional point in the stratified setting is that this product neighborhood must be regarded as a conically smooth stratified chart; once that is observed, the usual collar-gluing argument carries over without essential change.
\end{remark}

The next definition isolates the ordinary bordism category that the interval Segal space constructed later is meant to recover, after passage to the homotopy category. In particular, the choices of objects, morphisms, and composition below are designed so as to make that comparison precise.

\begin{definition}\label{ch5:: def: strat bord ord}
    The category $\cS\mathsf{Cob}_d$ of unoriented (conically smooth) stratified bordisms has:
    \begin{itemize}
        \item \textbf{Objects:} Objects are stratified closed $(d-1)$-manifolds.\footnote{By a closed stratified $d$-manifold we mean a stratified $d$-manifold whose underlying topological $d$-manifold has empty boundary and is compact.} 
        \item \textbf{Morphisms:} A morphism $M\maps \Sigma_{\sf in} \ra \Sigma_{\sf out}$ is an isomorphism class of $d$-dimensional stratified bordisms 
        \begin{equation}
            \left[\big(M, \varphi_{\sf in}, \varphi_{\sf out}\big)\right].
        \end{equation}
        \item \textbf{Composition:} Consider a pair of composable morphisms
        \begin{equation}
            \left[\big(M \maps \Sigma_0 \ra \Sigma_1,\varphi_{\sf in},\varphi_{\sf out}\big)\right], \qquad \left[\big(N \maps \Sigma_1 \ra \Sigma_2,\psi_{\sf in},\psi_{\sf out}\big)\right].
        \end{equation}
        Their composite is defined by gluing along the identified copies of $\Sigma_1$:
        \begin{equation}
            \left[\big(N,\psi_{\sf in},\psi_{\sf out}\big)\right]\circ \left[\big(M,\varphi_{\sf in},\varphi_{\sf out}\big)\right]
            \,\, \overset{\sf def}{=} \,\,
            \left[\big(M\cup_{\Sigma_1} N, q_M\circ \varphi_{\sf in}, q_N\circ \psi_{\sf out}\big)\right],
        \end{equation}
        where the gluing is performed along
        \begin{equation}
            \varphi_{\sf out}|_{\{0\}\times \Sigma_1}\maps \Sigma_1 \hookra M, \qquad \psi_{\sf in}|_{\{0\}\times \Sigma_1}\maps \Sigma_1 \hookra N,
        \end{equation}
        and $q_M$ and $q_N$ denote the canonical maps to the glued space.
        \item \textbf{Identity morphisms:} The identity on an object $\Sigma$ is represented by the cylinder bordism
        \begin{equation}
            \left[\big(\Sigma\times [0,1],\varphi_{\sf in},\varphi_{\sf out}\big)\right]\maps \Sigma \ra \Sigma,
        \end{equation}
        where $\varphi_{\sf in}$ and $\varphi_{\sf out}$ are the evident collars.
    \end{itemize}
\end{definition}

\begin{proposition}
    The above data defines a category. Moreover, $\cS\mathsf{Cob}_d$ is symmetric monoidal under disjoint union.
\end{proposition}
\begin{proof}
    The transitivity argument above shows that composition of composable morphisms again produces a stratified bordism. If
    \begin{equation}
        F\maps (M,\varphi_{\sf in},\varphi_{\sf out}) \xra{\;\cong\;} (M',\varphi'_{\sf in},\varphi'_{\sf out}), \qquad G\maps (N,\psi_{\sf in},\psi_{\sf out}) \xra{\;\cong\;} (N',\psi'_{\sf in},\psi'_{\sf out})
    \end{equation}
    are isomorphisms of composable stratified bordisms, then $F$ and $G$ glue to an isomorphism
    \begin{equation}
        M\cup_{\Sigma_1}N \xra{\;\cong\;} M'\cup_{\Sigma_1}N',
    \end{equation}
    so composition descends to isomorphism classes. The cylinder bordisms act as identities. Associativity of composition and the unit laws hold because collar-gluing is associative up to canonical isomorphism, and therefore strictly after passing to isomorphism classes. The symmetric monoidal structure is given on objects and morphisms by disjoint union.
\end{proof}

\section{An interval Segal space of \texorpdfstring{$d$}{d}-dimensional stratified bordisms}

We finally turn our attention to the construction of an interval Segal space of stratified bordisms. The approach is the same as in the smooth unstratified case: using intervals, together with appropriate conically smooth submersivity conditions, to construct a functor from $\Int$ landing in $\shv(\man)$, which, after post-composition with the concordance resolution, yields the desired stratified bordism $(\oo,1)$-category.

Additionally, we show that this stratified bordism $(\oo,1)$-category can be endowed with a symmetric monoidal structure and that its homotopy category is symmetric monoidal equivalent to \Cref{ch5:: def: strat bord ord}.



\begin{remark}
    We will want to consider stratified $d$-manifolds embedded in a non-stratified ambient space, namely a smooth manifold. In the case of interest this ambient space will be $\RR^{N}$, or more generally a finite-dimensional real vector space $V$. This requires the datum of a conically smooth stratified space $\big(M \ra P\big)$ together with a functor $\varphi \maps \exit(M) \ra \Vect^{\inj}_{|\sB\sO(d)}$ trivializing $\sT_{M}$ and a conically smooth embedding 
    \begin{equation}
        i \maps (M \ra P) \hookra X
    \end{equation}
    where $X$ is a smooth manifold, possibly with boundary, in which case one imposes the expected compatibility conditions along the boundary.
\end{remark}

\begin{definition}\label{ch5:: def: bord_d strat}
    Fix a finite-dimensional real vector space $V$. We define a functor 
    \begin{equation}
    \tilbordstr_{d,V}\maps \man^{\op}\times \Int \longrightarrow \set    
    \end{equation}
    sending a pair $(M,T)$, where $M\in \man^{\op}$ and $T\in \Int$, to the set of bounded closed subspaces 
    \begin{equation}
   \big( W\subseteq M\times \mathsf{cor}(T) \times V \big)
    \end{equation} 
    such that the following three conditions are satisfied:
    \begin{enumerate}
    \item \label{id: st 1} $W$ is endowed with the structure of a conically smooth stratification and a $\SMfd_{d}^{\partial}$-framing. In particular, it is a conically smooth stratified space with a boundary structure
    \begin{equation}
         (W \ra P),\quad (W \xra{c} \cP(*));
    \end{equation} 
    and the embedding must preserve the underlying stratifications, i.e., send the boundary of $W$ to the boundary of $M\times \core(T) \times V$.
    \item \label{id: st 2} The restricted projection 
    \begin{equation}
        \pi\maps \big(W \cap (M\times T^{\mathsf c} \times V) \big) \hookra  M \times T^{\mathsf c} \times V \fib  M \times T^{\mathsf c}
    \end{equation} 
    is a conically smooth submersion.
    \item \label{id: st 3} The restricted projection 
    \begin{equation}
    \pi_W\maps W \hookra M\times \mathsf{cor}(T) \times V \fib M 
    \end{equation}
    is a conically smooth submersion.
    \end{enumerate} 

    \noindent Now we define how $\tilbordstr_{d,V}$ acts on morphisms:
    \begin{itemize}
    \item[$(\man)$] The morphism $(f,\mathsf{id}_{T})\maps (M,T) \longra (N,T)$ maps the closed submanifold 
    \begin{equation}
    W\subseteq M\times \mathsf{cor}(T) \times V
    \end{equation}
    to the pullback along $(f\times \id)$:
    \begin{equation}
        \begin{tikzcd}[ampersand replacement=\&]
            {(f\times \id)^\ast W} \& {W} \\
            \\
            N\times \mathsf{cor}(T) \times V \& M\times \mathsf{cor}(T) \times V
            \arrow[from=1-1, to=1-2]
            \arrow[hook, "(f\times \id)^\ast \pi_{|W}"',from=1-1, to=3-1]
            \arrow["\lrcorner"{anchor=center, pos=0.125}, draw=none, from=1-1, to=3-2]
            \arrow[hook, from=1-2, to=3-2]
            \arrow["f\times \id"', from=3-1, to=3-2]
        \end{tikzcd}
    \end{equation}
    which is a conically smooth stratified $\big(d + \mathsf{dim}(N)\big)$-dimensional submanifold $f^\ast W \subseteq N \times \mathsf{cor}(T) \times V$.
    \item[$(\Int)$] A morphism $(\id_M, T \longra S)$ acts by \emph{forgetting} connected components and restricting the stratified $(d+\mathsf{dim}(M))$-manifold $(W\ra P)$.
    \end{itemize}  
\end{definition}

\begin{remark}
    This definition depends on the existence of a conically smooth embedding for every isomorphism class of compact stratified $d$-manifold with boundary. This is ensured by \Cref{ch4:: prop: neat embedding strat}.  
\end{remark}

\begin{observation}
    For every linear embedding of finite-dimensional real vector spaces $i\maps V \hookra W$, there is an induced natural transformation of bifunctors
    \begin{equation}
        \tilbordstr_{d,V} \xhookrightarrow{\,\,\,\,\,\,\,\,} \tilbordstr_{d,W}
    \end{equation}
    defined by post-composing the ambient embedding with $i$. We therefore define
    \begin{equation}
        \tilbordstr_{d} \,\, \overset{\sf def}{=} \,\, \hoColimit[V \in \vect^{\sf inj}]\,\tilbordstr_{d,V}
    \end{equation}
    as a bifunctor on $\man^{\op}\times \Int$. Applying the concordance resolution then gives
    \begin{equation}
        \bordstr_{d}(I) \,\, \overset{\sf def}{=} \,\, \big|\,\tilbordstr_{d}(\AA^\bullet,I)\,\big|
    \end{equation}
    for each $I\in \Int$. In this way we obtain a functor
    \begin{align}
        \bordstr_{d} \maps \Int &\longra \sset \\
        I &\longmapsto \bordstr_d(I).
    \end{align}
 
\end{observation}

The following is a critical observation rendering the definition of $\tilbordstr_{d,V}$ a sensible one.
\begin{observation}
    An important point for the soundness of the definition of $\tilbordstr_{d,V}$ is that if we pull back a conically smooth embedded submanifold, together with its conically smooth proper submersion over $M$, along a smooth map $f\maps M \longra N$, we again obtain a conically smooth proper submersion over $N$.
    
    From \Cref{ch4:: thm: cbl pullback}, we know that constructible bundles pull back along any conically smooth map. In particular, constructible bundles over a smooth base $K$ amount to conically smooth fiber bundles over $K$. From the stratified analogue of the Ehresmann fibration theorem, namely \Cref{ch4:: cor: stratified ehresmann boundary}, we know that 
    \begin{enumerate}
        \item[$(*)$] a proper conically smooth submersion $X \ra K$ over a smooth manifold $K$ is a conically smooth fiber bundle, 
    \end{enumerate} 
    ensuring stability under pullback for proper conically smooth submersions (over a smooth base $K$).
\end{observation}

\medskip

\begin{remark}
    To make contact with bordisms, let us depict an example and a pair of non-examples of a $0$-simplex in $\tilbordstr_2(\AA^{\bullet}, \sU(I))$, which illustrate how the construction captures objects in a stratified bordism category. The first diagram depicts a cylinder homeomorphic to $S^1\times [0,1]$ stratified in such a way that the tangent spaces to the $1$-dimensional stratum always project positively to the interval $I$, thereby yielding a valid $0$-simplex:
    \begin{equation}
        \tikzfig{figures/Stcob/SCO12}{\begin{tikzpicture}
	\begin{pgfonlayer}{nodelayer}
		\node [style=none] (0) at (-2, 3) {};
		\node [style=none] (1) at (-2, 1) {};
		\node [style=none] (2) at (3, 4.75) {};
		\node [style=none] (3) at (3, 2.75) {};
		\node [style=none] (4) at (-5, -1) {};
		\node [style=none] (5) at (7, -1) {};
		\node [style=none] (6) at (-2, -1) {};
		\node [style=none] (7) at (3, -1) {};
		\node [style=none] (8) at (-7.75, -3) {};
		\node [style=none] (9) at (4.25, -3) {};
		\node [style=none] (10) at (-4.75, -3) {};
		\node [style=none] (11) at (0.25, -3) {};
		\node [style=none] (36) at (-5, -1) {};
		\node [style=none] (37) at (-2, -1) {};
		\node [style=none] (38) at (-7.75, -3) {};
		\node [style=none] (39) at (-4.75, -3) {};
		\node [style=endpoint] (40) at (-5, -1) {};
		\node [style=endpoint] (41) at (-7.75, -3) {};
		\node [style=endpoint] (42) at (4.25, -3) {};
		\node [style=endpoint] (43) at (7, -1) {};
		\node [style=none] (44) at (3, -1) {};
		\node [style=none] (45) at (7, -1) {};
		\node [style=none] (46) at (0.25, -3) {};
		\node [style=none] (47) at (4.25, -3) {};
		\node [style=none] (48) at (3, 4.75) {};
		\node [style=none] (49) at (3, 2.75) {};
		\node [style=none] (50) at (-2, 1) {};
		\node [style=none] (51) at (-2, 3) {};
		\node [style=none] (52) at (-2.5, 2.25) {};
		\node [style=none] (53) at (2.55, 3.75) {};
		\node [style=none] (54) at (-2.4, 1.5) {};
		\node [style=none] (55) at (-1.6, 2.5) {};
		\node [style=0 strat red] (57) at (-2.375, 1.5) {};
		\node [style=0 strat red] (58) at (2.55, 3.75) {};
		\node [style=0 strat red] (59) at (-1.575, 2.5) {};
		\node [style=0 strat red] (60) at (3.45, 4) {};
		\node [style=none] (61) at (0.25, 3.5) {};
		\node [style=none] (62) at (1.5, 4.175) {};
	\end{pgfonlayer}
	\begin{pgfonlayer}{edgelayer}
		\draw [style=geometric, in=195, out=0] (0.center) to (2.center);
		\draw [style=geometric, in=195, out=-15, looseness=1.25] (1.center) to (3.center);
		\draw [style=geometric] (4.center) to (5.center);
		\draw [style={open_int}] (4.center) to (6.center);
		\draw [style={open_int}] (7.center) to (5.center);
		\draw [style=geometric] (8.center) to (9.center);
		\draw [style={open_int}] (8.center) to (10.center);
		\draw [style={open_int}] (11.center) to (9.center);
		\draw [style=defect line 3 dashed] (1.center) to (6.center);
		\draw [style=ts-dashed-red] (3.center) to (7.center);
		\draw [style=red shade low op] (36.center)
			 to (38.center)
			 to (39.center)
			 to (37.center)
			 to cycle;
		\draw [style=red shade low op] (47.center)
			 to (46.center)
			 to (44.center)
			 to (45.center)
			 to cycle;
		\draw [style=geometric, in=-180, out=180, looseness=0.75] (48.center) to (49.center);
		\draw [style=geometric, in=0, out=0, looseness=0.75] (48.center) to (49.center);
		\draw [style=geometric, in=180, out=-180, looseness=0.75] (50.center) to (51.center);
		\draw [style=ts-dashed, in=0, out=0, looseness=0.75] (51.center) to (50.center);
		\draw [style=defect line 3, in=-180, out=0, looseness=0.75] (54.center) to (53.center);
		\draw [style=defect line 3 ds, in=-120, out=30, looseness=1.25] (59) to (61.center);
		\draw [style=defect line 3, in=-105, out=-15] (61.center) to (62.center);
		\draw [style=defect line 3 ds, in=195, out=15] (62.center) to (60);
	\end{pgfonlayer}
\end{tikzpicture}
}
    \end{equation}
    \medskip

    \nid Next, we depict again a surface homeomorphic to $S^1\times [0,1]$, but whose 1-dimensional stratum admits two critical points relative to the projection to $I$, thereby failing to satisfy condition (\ref{id: st 2}) — and (\ref{id: st 3}), by extension:
    \begin{equation}
        \tikzfig{figures/Stcob/SCO13}{\begin{tikzpicture}
	\begin{pgfonlayer}{nodelayer}
		\node [style=none] (0) at (-2, 3) {};
		\node [style=none] (1) at (-2, 1) {};
		\node [style=none] (2) at (3, 4.75) {};
		\node [style=none] (3) at (3, 2.75) {};
		\node [style=none] (4) at (-5, -1) {};
		\node [style=none] (5) at (7, -1) {};
		\node [style=none] (6) at (-2, -1) {};
		\node [style=none] (7) at (3, -1) {};
		\node [style=none] (8) at (-7.75, -3) {};
		\node [style=none] (9) at (4.25, -3) {};
		\node [style=none] (10) at (-4.75, -3) {};
		\node [style=none] (11) at (0.25, -3) {};
		\node [style=none] (36) at (-5, -1) {};
		\node [style=none] (37) at (-2, -1) {};
		\node [style=none] (38) at (-7.75, -3) {};
		\node [style=none] (39) at (-4.75, -3) {};
		\node [style=endpoint] (40) at (-5, -1) {};
		\node [style=endpoint] (41) at (-7.75, -3) {};
		\node [style=endpoint] (42) at (4.25, -3) {};
		\node [style=endpoint] (43) at (7, -1) {};
		\node [style=none] (44) at (3, -1) {};
		\node [style=none] (45) at (7, -1) {};
		\node [style=none] (46) at (0.25, -3) {};
		\node [style=none] (47) at (4.25, -3) {};
		\node [style=none] (48) at (3, 4.75) {};
		\node [style=none] (49) at (3, 2.75) {};
		\node [style=none] (50) at (-2, 1) {};
		\node [style=none] (51) at (-2, 3) {};
		\node [style=none] (52) at (-2.5, 2.25) {};
		\node [style=none] (53) at (2.55, 3.75) {};
		\node [style=none] (54) at (-2.4, 1.5) {};
		\node [style=none] (55) at (-1.6, 2.5) {};
		\node [style=0 strat red] (57) at (-2.375, 1.5) {};
		\node [style=0 strat red] (58) at (2.55, 3.75) {};
		\node [style=0 strat red] (59) at (-1.575, 2.5) {};
		\node [style=0 strat red] (60) at (3.45, 4) {};
		\node [style=none] (61) at (-0.75, 3.15) {};
		\node [style=none] (62) at (1.5, 4.175) {};
		\node [style=none] (63) at (-0.25, 2.75) {};
	\end{pgfonlayer}
	\begin{pgfonlayer}{edgelayer}
		\draw [style=geometric, in=195, out=0] (0.center) to (2.center);
		\draw [style=geometric, in=195, out=-15, looseness=1.25] (1.center) to (3.center);
		\draw [style=geometric] (4.center) to (5.center);
		\draw [style={open_int}] (4.center) to (6.center);
		\draw [style={open_int}] (7.center) to (5.center);
		\draw [style=geometric] (8.center) to (9.center);
		\draw [style={open_int}] (8.center) to (10.center);
		\draw [style={open_int}] (11.center) to (9.center);
		\draw [style=defect line 3 dashed] (1.center) to (6.center);
		\draw [style=ts-dashed-red] (3.center) to (7.center);
		\draw [style=red shade low op] (36.center)
			 to (38.center)
			 to (39.center)
			 to (37.center)
			 to cycle;
		\draw [style=red shade low op] (47.center)
			 to (46.center)
			 to (44.center)
			 to (45.center)
			 to cycle;
		\draw [style=geometric, in=-180, out=180, looseness=0.75] (48.center) to (49.center);
		\draw [style=geometric, in=0, out=0, looseness=0.75] (48.center) to (49.center);
		\draw [style=geometric, in=180, out=-180, looseness=0.75] (50.center) to (51.center);
		\draw [style=ts-dashed, in=0, out=0, looseness=0.75] (51.center) to (50.center);
		\draw [style=defect line 3, in=270, out=0] (57) to (63.center);
		\draw [style=defect line 3, in=390, out=90] (63.center) to (61.center);
		\draw [style=defect line 3 ds, in=15, out=-150] (61.center) to (59);
		\draw [style=defect line 3 ds, in=15, out=-165] (60) to (62.center);
		\draw [style=defect line 3, in=-165, out=-120, looseness=1.25] (62.center) to (58);
	\end{pgfonlayer}
\end{tikzpicture}
}
    \end{equation} 
    \medskip

    \nid Lastly, we depict the same cylinder with a stratification containing a 0-dimensional stratum. At this point, the map fails to be a submersion for dimensionality reasons, thereby also failing to be a valid 0-simplex:
    \begin{equation}
        \tikzfig{figures/Stcob/SCO14}{\begin{tikzpicture}
	\begin{pgfonlayer}{nodelayer}
		\node [style=none] (0) at (-2, 3) {};
		\node [style=none] (1) at (-2, 1) {};
		\node [style=none] (2) at (3, 4.75) {};
		\node [style=none] (3) at (3, 2.75) {};
		\node [style=none] (4) at (-5, -1) {};
		\node [style=none] (5) at (7, -1) {};
		\node [style=none] (6) at (-2, -1) {};
		\node [style=none] (7) at (3, -1) {};
		\node [style=none] (8) at (-7.75, -3) {};
		\node [style=none] (9) at (4.25, -3) {};
		\node [style=none] (10) at (-4.75, -3) {};
		\node [style=none] (11) at (0.25, -3) {};
		\node [style=none] (36) at (-5, -1) {};
		\node [style=none] (37) at (-2, -1) {};
		\node [style=none] (38) at (-7.75, -3) {};
		\node [style=none] (39) at (-4.75, -3) {};
		\node [style=endpoint] (40) at (-5, -1) {};
		\node [style=endpoint] (41) at (-7.75, -3) {};
		\node [style=endpoint] (42) at (4.25, -3) {};
		\node [style=endpoint] (43) at (7, -1) {};
		\node [style=none] (44) at (3, -1) {};
		\node [style=none] (45) at (7, -1) {};
		\node [style=none] (46) at (0.25, -3) {};
		\node [style=none] (47) at (4.25, -3) {};
		\node [style=none] (48) at (3, 4.75) {};
		\node [style=none] (49) at (3, 2.75) {};
		\node [style=none] (50) at (-2, 1) {};
		\node [style=none] (51) at (-2, 3) {};
		\node [style=none] (52) at (-2.5, 2.25) {};
		\node [style=none] (53) at (2.55, 3.75) {};
		\node [style=none] (54) at (-2.4, 1.5) {};
		\node [style=none] (55) at (-1.6, 2.5) {};
		\node [style=0 strat red] (57) at (-2.375, 1.5) {};
		\node [style=0 strat red] (58) at (2.55, 3.75) {};
		\node [style=0 strat red] (59) at (-1.575, 2.5) {};
		\node [style=none] (61) at (-0.75, 3.15) {};
		\node [style=none] (62) at (1.5, 4.175) {};
		\node [style=none] (63) at (-0.25, 2.75) {};
		\node [style=0 strat red] (64) at (-0.25, 2.75) {};
	\end{pgfonlayer}
	\begin{pgfonlayer}{edgelayer}
		\draw [style=geometric, in=195, out=0] (0.center) to (2.center);
		\draw [style=geometric, in=195, out=-15, looseness=1.25] (1.center) to (3.center);
		\draw [style=geometric] (4.center) to (5.center);
		\draw [style={open_int}] (4.center) to (6.center);
		\draw [style={open_int}] (7.center) to (5.center);
		\draw [style=geometric] (8.center) to (9.center);
		\draw [style={open_int}] (8.center) to (10.center);
		\draw [style={open_int}] (11.center) to (9.center);
		\draw [style=defect line 3 dashed] (1.center) to (6.center);
		\draw [style=ts-dashed-red] (3.center) to (7.center);
		\draw [style=red shade low op] (36.center)
			 to (38.center)
			 to (39.center)
			 to (37.center)
			 to cycle;
		\draw [style=red shade low op] (47.center)
			 to (46.center)
			 to (44.center)
			 to (45.center)
			 to cycle;
		\draw [style=geometric, in=-180, out=180, looseness=0.75] (48.center) to (49.center);
		\draw [style=geometric, in=0, out=0, looseness=0.75] (48.center) to (49.center);
		\draw [style=geometric, in=180, out=-180, looseness=0.75] (50.center) to (51.center);
		\draw [style=ts-dashed, in=0, out=0, looseness=0.75] (51.center) to (50.center);
		\draw [style=defect line 3, in=270, out=0] (57) to (63.center);
		\draw [style=defect line 3, in=375, out=90] (63.center) to (61.center);
		\draw [style=defect line 3 ds] (61.center) to (59);
		\draw [style=defect line 3, in=-165, out=15] (64) to (58);
	\end{pgfonlayer}
\end{tikzpicture}
}.
    \end{equation}
    \medskip

    \nid These examples highlight the fact that the second and third examples do not represent objects in a $2$-dimensional stratified bordism category, whereas the first does. Up to reparametrization, it captures the circle with two $0$-dimensional strata
    \begin{equation}
        \tikzfig{figures/Stcob/SCO15}{\begin{tikzpicture}
	\begin{pgfonlayer}{nodelayer}
		\node [style=none] (3) at (0, -1.5) {};
		\node [style=none] (48) at (0, 1.5) {};
		\node [style=none] (49) at (0, -1.5) {};
		\node [style=0 strat red] (58) at (-1.525, 0) {};
		\node [style=0 strat red] (60) at (1.425, 0.575) {};
	\end{pgfonlayer}
	\begin{pgfonlayer}{edgelayer}
		\draw [style=geometric, in=-180, out=180, looseness=1.75] (48.center) to (49.center);
		\draw [style=geometric, in=0, out=0, looseness=1.75] (48.center) to (49.center);
	\end{pgfonlayer}
\end{tikzpicture}
}.
    \end{equation}
    \medskip

    To get a sense of how bordisms are encoded, let us pick out an example of a $0$-simplex in $\tilbordstr_2(\AA^\bullet,\sU(I,J))$. Over the red-shaded region in the middle there are different singular points relative to the interval projection: the $0$-dimensional stratum in red and the saddle point of the pair-of-pants in the $1$-dimensional stratum.

    \begin{equation}
        \tikzfig{figures/Stcob/SCO16}{\begin{tikzpicture}
	\begin{pgfonlayer}{nodelayer}
		\node [style=none] (0) at (-2, 3) {};
		\node [style=none] (1) at (-2, 1) {};
		\node [style=none] (2) at (6.25, 3.75) {};
		\node [style=none] (3) at (6.25, 1.75) {};
		\node [style=none] (4) at (-5, -1) {};
		\node [style=none] (5) at (10, -1) {};
		\node [style=none] (6) at (-2, -1) {};
		\node [style=none] (7) at (3, -1) {};
		\node [style=none] (8) at (-7.75, -3) {};
		\node [style=none] (9) at (7.25, -3) {};
		\node [style=none] (10) at (-4.75, -3) {};
		\node [style=none] (11) at (0.25, -3) {};
		\node [style=none] (36) at (-5, -1) {};
		\node [style=none] (37) at (-2, -1) {};
		\node [style=none] (38) at (-7.75, -3) {};
		\node [style=none] (39) at (-4.75, -3) {};
		\node [style=endpoint] (40) at (-5, -1) {};
		\node [style=endpoint] (41) at (-7.75, -3) {};
		\node [style=endpoint] (42) at (7.25, -3) {};
		\node [style=endpoint] (43) at (10, -1) {};
		\node [style=none] (44) at (3, -1) {};
		\node [style=none] (45) at (10, -1) {};
		\node [style=none] (46) at (0.25, -3) {};
		\node [style=none] (47) at (7.25, -3) {};
		\node [style=none] (48) at (6.25, 3.75) {};
		\node [style=none] (49) at (6.25, 1.75) {};
		\node [style=none] (50) at (-2, 1) {};
		\node [style=none] (51) at (-2, 3) {};
		\node [style=none] (52) at (-2.5, 2.25) {};
		\node [style=none] (54) at (-2.4, 1.5) {};
		\node [style=none] (55) at (-1.6, 2.5) {};
		\node [style=0 strat red] (57) at (-2.375, 1.5) {};
		\node [style=0 strat red] (59) at (-1.575, 2.5) {};
		\node [style=none] (61) at (4.5, 3.35) {};
		\node [style=none] (63) at (5, 2.75) {};
		\node [style=0 strat red] (64) at (5, 2.75) {};
		\node [style=none] (75) at (3, -1) {};
		\node [style=none] (76) at (17.75, -1) {};
		\node [style=none] (77) at (9, -1) {};
		\node [style=none] (78) at (13.75, -1) {};
		\node [style=none] (79) at (7.25, -3) {};
		\node [style=none] (80) at (15, -3) {};
		\node [style=none] (81) at (3.25, -3) {};
		\node [style=none] (82) at (11, -3) {};
		\node [style=none] (84) at (9, -1) {};
		\node [style=none] (85) at (0.25, -3) {};
		\node [style=none] (86) at (3.25, -3) {};
		\node [style=endpoint] (89) at (15, -3) {};
		\node [style=endpoint] (90) at (17.75, -1) {};
		\node [style=none] (91) at (13.75, -1) {};
		\node [style=none] (92) at (17.75, -1) {};
		\node [style=none] (93) at (11, -3) {};
		\node [style=none] (94) at (15, -3) {};
		\node [style=none] (95) at (8.5, 5) {};
		\node [style=none] (96) at (8.25, 1) {};
		\node [style=none] (97) at (9.75, 3.5) {};
		\node [style=none] (98) at (9.5, 2) {};
		\node [style=none] (99) at (8.75, 2.75) {};
		\node [style=none] (100) at (13.75, 0.25) {};
		\node [style=none] (101) at (13.75, 1.25) {};
		\node [style=none] (102) at (13.75, 3.75) {};
		\node [style=none] (103) at (13.75, 5) {};
		\node [style=none] (104) at (13.75, 3.75) {};
		\node [style=none] (105) at (13.75, 5) {};
		\node [style=0 strat red] (106) at (13.5, 4.5) {};
		\node [style=0 strat red] (107) at (13.45, 0.75) {};
		\node [style=none] (108) at (7.5, 2) {};
		\node [style=none] (109) at (13.75, -1) {};
	\end{pgfonlayer}
	\begin{pgfonlayer}{edgelayer}
		\draw [style=geometric, in=195, out=0] (0.center) to (2.center);
		\draw [style=geometric, in=195, out=-15, looseness=1.25] (1.center) to (3.center);
		\draw [style=geometric] (4.center) to (5.center);
		\draw [style={open_int}] (4.center) to (6.center);
		\draw [style=geometric] (8.center) to (9.center);
		\draw [style={open_int}] (8.center) to (10.center);
		\draw [style={open_int}] (11.center) to (9.center);
		\draw [style=defect line 3 dashed] (1.center) to (6.center);
		\draw [style=red shade low op] (36.center)
			 to (38.center)
			 to (39.center)
			 to (37.center)
			 to cycle;
		\draw [style=red shade low op] (47.center)
			 to (46.center)
			 to (44.center)
			 to (45.center)
			 to cycle;
		\draw [style=geometric, in=180, out=-180, looseness=0.75] (50.center) to (51.center);
		\draw [style=ts-dashed, in=0, out=0, looseness=0.75] (51.center) to (50.center);
		\draw [style=defect line 3, in=240, out=0] (57) to (63.center);
		\draw [style=defect line 3, in=375, out=90] (63.center) to (61.center);
		\draw [style=defect line 3 ds, bend left=15] (61.center) to (59);
		\draw [style=geometric] (75.center) to (76.center);
		\draw [style={open_int}] (78.center) to (76.center);
		\draw [style=geometric] (79.center) to (80.center);
		\draw [style={open_int}] (82.center) to (80.center);
		\draw [style=red shade low op] (94.center)
			 to (93.center)
			 to (91.center)
			 to (92.center)
			 to cycle;
		\draw [style={open_int}] (7.center) to (5.center);
		\draw [style=geometric, in=-165, out=15] (48.center) to (95.center);
		\draw [style=geometric, in=165, out=0] (49.center) to (96.center);
		\draw [style=geometric, in=-90, out=150] (98.center) to (99.center);
		\draw [style=geometric, in=-180, out=90] (99.center) to (97.center);
		\draw [style=geometric, in=180, out=-15] (96.center) to (100.center);
		\draw [style=geometric, in=-180, out=-15, looseness=0.50] (98.center) to (101.center);
		\draw [style=geometric] (97.center) to (102.center);
		\draw [style=geometric, in=180, out=15, looseness=0.25] (95.center) to (103.center);
		\draw [style=geometric, in=-180, out=-180] (101.center) to (100.center);
		\draw [style=geometric, in=-180, out=180, looseness=0.75] (105.center) to (104.center);
		\draw [style=geometric, in=0, out=0, looseness=0.75] (101.center) to (100.center);
		\draw [style=geometric, in=0, out=0, looseness=0.75] (105.center) to (104.center);
		\draw [style=defect line 3, in=-165, out=0, looseness=0.75] (64) to (49.center);
		\draw [style=defect line 3 dashed, in=-165, out=0, looseness=0.50] (49.center) to (95.center);
		\draw [style=defect line 3, in=180, out=0] (95.center) to (106);
		\draw [style=defect line 3, in=135, out=30] (64) to (108.center);
		\draw [style=defect line 3, in=180, out=-45, looseness=0.75] (108.center) to (107);
		\draw [style=defect line 3 dashed] (100.center) to (109.center);
	\end{pgfonlayer}
\end{tikzpicture}
}
    \end{equation}
    \medskip 
\end{remark}

\begin{observation}
    The structure and intent behind \Cref{ch5:: def: bord_d strat} are manifestly analogous to those of \Cref{ch3:: def: bord_d} and \Cref{ch3:: def: bord_d tangential}. In particular, our path toward the interval-Segal model of the stratified bordism category follows the same sequence of steps: we first establish the analogous locality statements (\Cref{ch3:: prop: naive sheaf con} and \Cref{ch3:: prop: sheaf condition}) for $\tilbordstr_d$, and then prove the interval Segal condition.
\end{observation}

Just as in the unstratified smooth case, passing to the adjoint via the mapping-power adjunction produces an $\oo$-sheaf on $\man$.
\begin{proposition}\label{ch5:: prop: sheaf condition}
    The adjoint functor 
    \begin{equation}
    \tilbordstr_{d,V}(-,-)\maps \Int \longra \Fun(\man^{\op},\sset)
    \end{equation}
    factors through $\shv(\man)$, the category of $\oo$-sheaves on $\man$ relative to the Grothendieck topology of open covers. In other words, for every $T\in \Int$ the functor
    \[
    \tilbordstr_{d,V}(-,T)\maps \man^{\op}\longra \set
    \]
    is a smooth $\oo$-sheaf of (discrete) simplicial sets.
\end{proposition}
\begin{proof}
    The proof is entirely analogous to that of \Cref{ch3:: prop: sheaf condition}. The key observation in the stratified case is that being a conically smooth submersion is also an open condition — ensuring gluing satisfies conditions (\ref{id: st 1})--(\ref{id: st 3}).
\end{proof}

\begin{proposition}[Naive sheaf condition in the interval direction]\label{ch5:: prop: naive sheaf con}
Let $T\in \Int$, and let $\{U,V\}$ be an open cover of $\core(T)\subset [0,1]$. For each $k\geq 0$, the canonical map
\begin{equation}
\label{ch5:: eq: naive sheaf interval levelwise}
\tilbordstr_{d,V}\big(\AA^{k},T\big)\longrightarrow
\tilbordstr_{d,V}\big(\AA^{k},T\big)_{|U}\underset{\tilbordstr_{d,V}\big(\AA^{k},T\big)_{|U\cap V}}{\times}
\tilbordstr_{d,V}\big(\AA^{k},T\big)_{|V}
\end{equation}
is a bijection of sets. Equivalently, the simplicial set $\tilbordstr_{d,V}(\AA^{\bullet},T)$ is the strict pullback
\begin{equation}
\label{ch5:: eq: naive sheaf interval simplicial}
\tilbordstr_{d,V}(\AA^{\bullet},T)\;\cong\;
\tilbordstr_{d,V}(\AA^{\bullet},T)_{|U}\underset{\tilbordstr_{d,V}(\AA^{\bullet},T)_{|U\cap V}}{\times}
\tilbordstr_{d,V}(\AA^{\bullet},T)_{|V},
\end{equation}
where
\begin{equation}
\label{ch5:: eq: bord-restriction-U}
\tilbordstr_{d,V}(\AA^{k},T)_{|U}
=
\Big\{
W\subseteq \AA^{k}\times\big(U\cap \core(T)\big)\times V
\Big\}.
\end{equation}
Here $\tilbordstr_{d,V}(\AA^{k},T)_{|U}$ denotes the set of all $(d+k)$-dimensional stratified submanifolds with boundary
\[
W\subseteq \AA^{k}\times\big(U\cap \core(T)\big)\times V
\]
which satisfy conditions (\ref{id: st 1}), (\ref{id: st 2}), and (\ref{id: st 3}) from \Cref{ch5:: def: bord_d strat}, and the restriction maps are given by intersecting with
\[
\AA^{k}\times\big(U\cap V\cap \core(T)\big)\times V\subseteq 
\AA^{k}\times\big(U\cap\core(T)\big)\times V,
\]
and similarly for $V$.
\end{proposition}

\begin{proof}
    This is the stratified analogue of \Cref{ch3:: prop: naive sheaf con}, and the same argument applies. Fix $k\geq 0$ and write
    \begin{equation}
    Z = \AA^{k}\times \core(T)\times V.
    \end{equation}
    Injectivity of \eqref{ch5:: eq: naive sheaf interval levelwise} is immediate, since an element
    \begin{equation}
    W\subseteq Z
    \end{equation}
    of $\tilbordstr_{d,V}(\AA^{k},T)$ is recovered from its restrictions to
    $\AA^{k}\times(U\cap\core(T))\times V$ and
    $\AA^{k}\times(V\cap\core(T))\times V$.

    For surjectivity, let
    \begin{equation}
    W_U\subseteq \AA^{k}\times\big(U\cap\core(T)\big)\times V,\qquad
    W_V\subseteq \AA^{k}\times\big(V\cap\core(T)\big)\times V
    \end{equation}
    be compatible elements, and define $W := W_U\cup W_V\subseteq Z$. Since the two pieces agree on the overlap, this is a well-defined subset. Exactly as in Chapter 3, all the relevant structure is local on the source. Hence $W$ inherits the structure of a conically smooth stratified manifold with boundary by gluing the local models carried by $W_U$ and $W_V$. The boundary condition in (\ref{id: st 1}) is local, and the submersivity conditions (\ref{id: st 2}) and (\ref{id: st 3}) hold on $W$ because conically smooth submersion is an open condition. Therefore $W$ lies in $\tilbordstr_{d,V}(\AA^{k},T)$ and restricts to $(W_U,W_V)$.

    This proves surjectivity, hence bijectivity of \eqref{ch5:: eq: naive sheaf interval levelwise}. The simplicial-set statement follows levelwise.
\end{proof}

These two ingredients, in their unstratified versions, are the fundamental keys to proving the interval Segal condition in Chapter 3. For that reason, we will now introduce the analogue of $\bord_d$ for stratified $d$-manifolds, using the concordance resolution, and prove that the resulting pre-interval space is in fact an interval Segal space.

\begin{theorem}\label{ch5:: thm: stratbord interval segal}
    Let $d \geq 0$ be a non-negative integer and let $V$ be a finite-dimensional vector space of sufficiently high dimension. The composite functor
    \begin{align}
        \bordstr_{d,V} \maps \Int &\longra \shv(\man)  \longra \sset  \\
         I &\longmapsto \tilbordstr_{d,V}(-,I) \longmapsto \bB\tilbordstr_{d,V}(*,I) = \big|\,\tilbordstr_{d,V}(\AA^\bullet,I)\,\big| 
    \end{align}
    is an interval Segal space.
\end{theorem}

\begin{proof}
    The proof is broken down into two parts: 1) showing that $\bordstr_{d,V}$ sends morphisms in $\cW_{\pi_{0}}$ to weak equivalences of simplicial sets and 2) checking the interval Segal condition, which states — in one of its equivalent forms — that for every $\sU(A,B,C)\in \Int$, the following is a homotopy pullback square
     \begin{equation}\label{ch5:: eq: Segal condition for bordism}
        \begin{tikzcd}[ampersand replacement=\&]
            {\bordstr_{d,V}\big(U(A,B,C)\big)} \&\& {\bordstr_{d,V}\big(U(A,B)\big)} \\
            \\
            {\bordstr_{d,V}\big(U(B,C)\big)} \&\& {\bordstr_{d,V}\big(U(B)\big).}
            \arrow[from=1-1, to=1-3]
            \arrow[from=1-1, to=3-1]
            \arrow["h"{description, pos=0.3}, "\lrcorner"{anchor=center, pos=0.125}, draw=none, from=1-1, to=3-3]
            \arrow[from=1-3, to=3-3]
            \arrow[from=3-1, to=3-3]
        \end{tikzcd}
    \end{equation} 
    For 1), given a morphism 
    \begin{equation}
        T \hookra S
    \end{equation}
    inducing an isomorphism on $\pi_0$, the same diffeomorphism $s$ which interpolates between the two intervals $T$ and $S$ also provides an explicit simplicial homotopy equivalence between $\bordstr_{d,V}(S)$ and $\bordstr_{d,V}(T)$. For additional details see the first part of the proof of \Cref{ch3:: thm: interval segal bordism cat}.

    For 2), the only ingredient of the proof in the unstratified case which is not abstractly categorical and depends on the geometric context is the \emph{naive sheaf condition}, which is easily adapted to the stratified case in \Cref{ch5:: prop: naive sheaf con}. Once again, the details of the second part of the proof of \Cref{ch3:: thm: interval segal bordism cat} are the same.
\end{proof}

\begin{observation}
    We chose to first present this result, together with a complete proof of its unstratified variant, in Chapter 3 precisely because it adapts so readily to more general contexts, such as the present stratified one. In fact, we expect that this method for building higher bordism-type categories should apply in still broader settings, for instance to graph bordisms.
\end{observation}

Through disjoint union of conically smooth stratified $d$-manifolds, we can endow $\bordstr_d$ with a symmetric monoidal structure. This structure is also built in an entirely analogous fashion to that of the unstratified case. We recall the construction here for completion.

\begin{construction}\label{ch5:: cons: Gamma-space}
  For every $m\geq 0$, let $\bordstr_d^{\langle m \rangle}$ denote the following interval Segal space:
\begin{align}
    \bordstr_d^{\langle m \rangle} \maps \Int &\longra \sset \\
    I &\longmapsto \Big(\,[k] \mapsto \left\{M \subset \AA^{k}\times \mathsf{cor}(I)\times \RR^{\oo}\right\}\,\Big)
\end{align} 
where $M$ is given by a disjoint union $M = M_{1} \sqcup \cdots \sqcup M_{m}$ where each $M_{i}\subset \AA^{k} \times \mathsf{cor}(I)\times \RR^{\oo}$ is a $k$-simplex of $\bordstr_d(I)$. Given a pointed map of pointed finite sets, $f\maps \langle m \rangle \longra \langle n \rangle$, we construct a map of interval Segal spaces. This is the natural transformation whose component on $I \in \Int$ is
\begin{equation}
    \begin{tikzcd}
        {\bordstr_d^{\langle m \rangle }(I)} & {M = M_1 \sqcup \cdots \sqcup M_m} \\
        \\
        {\bordstr_d^{\langle n \rangle }(I)} & {(f_*)_I = \Big(\coprod_{f(i) = 1}M_i\Big) \sqcup \cdots \sqcup \Big(\coprod_{f(i) = n}M_i\Big)}
        \arrow[from=1-1, to=3-1]
        \arrow["{{(f_*)_I}}"', maps to, from=1-2, to=3-2]
    \end{tikzcd}
\end{equation}
where we omit the ambient data of the embeddings for simplicity.
\end{construction}

\Cref{ch5:: cons: Gamma-space} can again be proven to induce a symmetric monoidal structure, in this case on $\bordstr_d$.

\begin{theorem}\label{ch5:: thm: symmetric monoidal stratified}
    The assignment 
    \begin{align}
        \bordstr_{d}^{\langle \bullet \rangle} \maps \Gamma &\longra \mathsf{ISSp} \\
        \langle m \rangle & \longmapsto \bordstr_d^{\langle m \rangle}
    \end{align}
    is functorial and endows the interval Segal space $\bordstr_d$ of unoriented $d$-dimensional stratified bordisms with a symmetric monoidal structure.
\end{theorem}

\begin{proof}
    The details of the proof follow those of \Cref{ch3:: thm: symmetric monoidal}. Both proofs rely essentially on the contractibility of the space of conically smooth embeddings of a conically smooth stratified space $M$ into $\RR^{\oo}$ (see \Cref{ch4:: prop: contractibility of embeddings}). 
\end{proof}

\begin{lemma}\label{ch5:: lem: normalized neat embedding}
    Let 
    \begin{equation}
        \big(M,\varphi_{\sf in},\varphi_{\sf out}\big)\maps \Sigma_{\sf in} \ra \Sigma_{\sf out}
    \end{equation}
    be a compact stratified bordism. After possibly shrinking the given collars, there exists a neat embedding
    \begin{equation}
        E\maps M \hookra [0,1]\times \RR^{m}
    \end{equation}
    for some $m\geq 0$, such that
    \begin{equation}
        E(\partial_{\sf in}M)\subseteq \{0\}\times \RR^{m}, \qquad E(\partial_{\sf out}M)\subseteq \{1\}\times \RR^{m},
    \end{equation}
    and, on smaller collar neighborhoods,
    \begin{equation}
        E\circ \varphi_{\sf in}(t,x)=\big(t,e_{\sf in}(x)\big), \qquad E\circ \varphi_{\sf out}(s,x)=\big(1+s,e_{\sf out}(x)\big),
    \end{equation}
    for suitable conically smooth embeddings $e_{\sf in}$ and $e_{\sf out}$ into $\RR^{m}$.
\end{lemma}
\begin{proof}
    This is a minor adaptation of \Cref{ch4:: prop: neat embedding strat}. By compactness, after shrinking collars, we may assume that for some $\epsilon>0$ the images of
    \begin{equation}
        \varphi_{\sf in}\big([0,2\epsilon)\times \Sigma_{\sf in}\big), \qquad \varphi_{\sf out}\big((-2\epsilon,0]\times \Sigma_{\sf out}\big)
    \end{equation}
    are disjoint. Choose a proper conically smooth embedding $e\maps M \hookra \RR^{m}$ by \Cref{ch4:: prop: weak Whit strat}. Using \Cref{lem:partitions of unity}, we construct a conically smooth function
    \begin{equation}
        h\maps M \longra [0,1]
    \end{equation}
    such that
    \begin{equation}
        h\circ \varphi_{\sf in}(t,x)=t, \qquad h\circ \varphi_{\sf out}(s,x)=1+s
    \end{equation}
    on the smaller collars $[0,\epsilon)\times \Sigma_{\sf in}$ and $(-\epsilon,0]\times \Sigma_{\sf out}$, and such that $h(M\setminus \partial M)\subset (0,1)$.

    As in the proof of \Cref{ch4:: prop: neat embedding strat}, after modifying $e$ by a conically smooth bump function near the boundary, we may assume that $e$ is constant along the collar direction on these smaller collar neighborhoods. Therefore the product
    \begin{equation}
        E=(h,e)\maps M \longra [0,1]\times \RR^{m}
    \end{equation}
    is a conically smooth embedding, and on the chosen collar neighborhoods it takes the form
    \begin{equation}
        E\circ \varphi_{\sf in}(t,x)=\big(t,e_{\sf in}(x)\big), \qquad E\circ \varphi_{\sf out}(s,x)=\big(1+s,e_{\sf out}(x)\big).
    \end{equation}
    In particular, $E$ is neat and carries the incoming and outgoing boundary components to $\{0\}\times \RR^{m}$ and $\{1\}\times \RR^{m}$, respectively.
\end{proof}

\section{Homotopy category as the stratified bordism 1-category}

    Before studying $(\oo,1)$-categories we studied ordinary 1-categories, and instances of bordism categories date back at least to Atiyah's 1988 paper \cite{At88}. To confirm that the interval Segal space $\bordstr_d$ really lifts the ordinary stratified bordism category, we should compute its symmetric monoidal homotopy $1$-category and compare it with the category $\cS\mathsf{Cob}_d$ constructed earlier in this chapter.

\begin{remark}\label{ch5:: rmk: homotopy ehresmann strat}
    In the next theorem we work with the homotopy category of the Segal space
    \begin{equation}
        X \,\, \overset{\sf def}{=} \,\, \mathsf{RKan}_{\gamma}(\bordstr_d).
    \end{equation}
    Its objects are therefore coherent families of compact conically smooth stratified $(d-1)$-manifolds. Exactly as in the smooth case, each such family determines a unique isomorphism class of stratified manifold, obtained by taking the fiber over any point in the relevant interval. The point is that the projection to the interval is a proper conically smooth submersion, so \Cref{ch4:: cor: stratified ehresmann boundary} identifies all of its fibers up to conically smooth stratified diffeomorphism. The same observation applies to 1-simplices, where one obtains a stratified bordism by restricting to the closed subinterval between the chosen incoming and outgoing midpoints.
\end{remark}

\begin{theorem}\label{ch5:: thm: homotopy category stratified bord}
    There is an equivalence of symmetric monoidal 1-categories
    \begin{equation}
        \ho(\bordstr_d) \,\, \simeq \,\, \cS\mathsf{Cob}_d
    \end{equation}
    between the homotopy category of $\bordstr_d$ and the ordinary category of unoriented stratified bordisms.
\end{theorem}

\begin{proof}
    The proof follows the same pattern as the corresponding homotopy-category theorem in Chapter 3, with the smooth geometric inputs replaced by their stratified analogues from Chapter 4.

    We define a functor
    \begin{equation}
        h\maps \ho(\bordstr_d) \longra \cS\mathsf{Cob}_d.
    \end{equation}
    On objects, if $(\Sigma_I)_{I\in \cJ_0}$ is a coherent family representing an object of $\ho(\bordstr_d)$, then we set
    \begin{equation}
        h\big((\Sigma_I)_{I\in \cJ_0}\big) \,\, \overset{\sf def}{=} \,\, \pi_I^{-1}\big(\mathsf{mid}(I)\big)
    \end{equation}
    for any choice of interval $I\in \cJ_0$. By \Cref{ch5:: rmk: homotopy ehresmann strat}, this is well-defined up to conically smooth stratified diffeomorphism. On morphisms, represented by coherent families of embedded stratified bordisms over some $\sU(A,B)\in \cJ_1$, we restrict to the closed subinterval between the chosen midpoints and take the resulting isomorphism class in $\cS\mathsf{Cob}_d$. The same stratified Ehresmann argument shows that this is independent of all auxiliary choices.

    Essential surjectivity and fullness are proved exactly as in Chapter 3, except that one uses \Cref{ch5:: lem: normalized neat embedding} in place of the smooth neat-embedding statement. Indeed, every compact stratified bordism admits a normalized neat embedding into $[0,1]\times \RR^m$, and after choosing for each $\sU(A,B)\in \cJ_1$ a diffeomorphism
    \begin{equation}
        [0,1] \xra{\,\,\,\simeq\,\,\,} \core\big(\sU(A,B)\big)
    \end{equation}
    carrying $0$ and $1$ to the prescribed incoming and outgoing midpoints, these embeddings assemble into the required coherent family. The contractibility of the relevant space of interval reparameterizations is the same as in the smooth case.

    For faithfulness, suppose two morphisms in $\ho(\bordstr_d)$ have the same image in $\cS\mathsf{Cob}_d$. Then the extracted stratified bordisms are isomorphic, so after choosing a representative compact stratified bordism
    \begin{equation}
        \big(M,\varphi_{\sf in},\varphi_{\sf out}\big)\maps \Sigma \ra \Sigma'
    \end{equation}
    for that common isomorphism class, both morphisms may be represented by normalized neat embeddings of $M$ into $[0,1] \times \RR^{\oo}$. After fixing the same collar-compatible height function $h \maps M \ra [0,1]$ as in \Cref{ch5:: lem: normalized neat embedding}, the difference between the two representatives lies only in the Euclidean embedding component. By the weak contractibility of the space of conically smooth embeddings into $\RR^{\oo}$ (\Cref{ch4:: prop: contractibility of embeddings}), these two embedding components lie in the same path component, hence determine concordant $1$-simplices in $\bordstr_{d}$. Therefore they define the same morphism in $\ho(\bordstr_{d})$, so $h$ is faithful.

    It follows that $h$ is an equivalence of categories. Compatibility with the symmetric monoidal structures follows from the fact that both sides are defined by disjoint union, exactly as in the smooth case discussed in Chapter 3.
\end{proof}

\begin{remark}
    Before moving on, it is worth recording a possible further direction suggested by the preceding construction. Even in the stratified setting developed in this chapter, the parameter direction still ranges over smooth manifolds. More precisely, the input is again a bifunctor
    \begin{equation}
        \Int \times \man^{\op} \longra \set,
    \end{equation}
    which is then transposed and post-composed with a concordance-resolution procedure in order to obtain a space-valued interval Segal object. In this way, the interval variable is responsible for bordism composition, while the second variable is used to encode the spatial direction. The stratified character of the theory is therefore encoded in the fibers, not in the parameter category itself.

    This suggests a different possible generalization. One could try to replace the smooth parameter direction by a genuinely stratified one, and begin instead with a bifunctor
    \begin{equation}
        \Int \times \strat^{\op}\longra \set,
    \end{equation}
    in which the role played above by conically smooth submersions over a smooth base is replaced by the requirement of being a constructible bundle over a stratified base. Transposing such a bifunctor would produce
    \begin{equation}
        \Int \longra \Fun(\strat^{\op},\set),
    \end{equation}
    and one may then ask whether applying the topologizing diagram levelwise yields an assignment
    \begin{equation}
        \Int \longra \Fun(\strat^{\op},\sset)
    \end{equation}
    which lands in striation sheaves. Since striation sheaves form a model for $(\oo,1)$-categories \cite{AFR1}, this would mean that each interval object determines an $(\oo,1)$-category rather than merely a space. If, moreover, the resulting assignment satisfied the interval Segal condition internally, one would obtain an interval Segal object valued in $(\oo,1)$-categories, hence a double $(\oo,1)$-category, or equivalently a model for an $(\oo,2)$-category. This point of view has begun making appearances via symmetric monoidal double higher bordism categories in the recent work of Yang Yang \cite{YY22} and more recently in work of Krannich--Kupers \cite{Ku22}.

    From this perspective, the natural ambient candidate for the vertical direction is provided by constructible bundles, hence by $\Bun$. The point is not that one must decide in advance which particular subcategory of $\Bun$ is the correct one. On the contrary, that flexibility seems to be part of the picture: the work of Ayala--Francis--Rozenblyum suggests that $\Bun$ is itself a rich ambient object, with natural specializations to embeddings, refinements, closed morphisms, and related classes of stratified maps \cite{AFR1,AFR2}. If such a program can be carried out, then the horizontal direction would still encode bordism composition, while the vertical direction would retain geometric operations that are invisible once one passes directly to spaces.
\end{remark}

\section{Tangential structures}


Tangential structures on conically smooth stratified spaces form a rich and still largely unexplored theory. In this thesis we do not attempt to develop such a theory in full generality. Instead, the goal of the present section is more modest and more targeted: we want to isolate a flexible piece of \emph{local} data on stratified manifolds that plays, for defect theories, the same organizational role that ordinary tangential structures (orientation, framing, $G$-structure, \dots) play in the classical bordism category. In particular, we want a framework in which strata of different depths can carry \emph{labels} (or ``colors'') in a way that satisfies a kind of descent relative to the stratification. 

The guiding principle is that, for stratified spaces, local data is controlled by \emph{basics} — especially their links. This suggests that defect label data should be specified locally on the category of basics and
transported functorially along stratified embeddings between them. Concretely, we model a discrete
system of defect labels as a right fibration over the $(\infty,1)$-category of basics, and we focus on
the case of \emph{discrete} fibers. Via straightening, such an object can be interpreted as a
set-valued presheaf on the category of basics: it assigns to each local model a set of admissible
labels, together with restriction maps along morphisms of basics. From this point of view, a labeling of
a stratified manifold is naturally expressed as a functor out of its enter-path category, and the
functoriality built into the presheaf guarantees compatibility with refinement of stratifications and with gluing.

Besides the last section of Chapter 4, which discusses tangential structures on conically smooth stratified spaces, we recommend the recent survey by Lukas M\"uller \cite{Mu25}, as well as Ayala--Francis--Rozenblyum's original article \cite{AFR2}.


\subsection{Discrete defect data}

\begin{definition}
    We call a right fibration 
    \begin{equation}
        \Bsc_{d}^{\cS_{d-1},\DD} \longra \Bsc_{d}^{\cS_{d-1}}
    \end{equation}
    with discrete fibers a $d$\define{-dimensional defect datum}.
\end{definition}

\nid To understand this definition, we will unpack it in increasing degrees of complexity. We begin with a simple low-dimensional example.

\begin{observation}
    Recall that one of the significant features of conically smooth stratified spaces is that the obvious inclusion
    \begin{equation}
        \Aut\big(\sC(X)\big) \, \xla{\,\,\,\simeq\,\,\,} \, \Aut(X)
    \end{equation}
    is a homotopy equivalence for every compact conically smooth stratified space. An important consequence of this result is a fully homotopical characterization of the spaces of morphisms in a category of basics:
    \begin{equation}
        \Strat^{\sf emb}\big(\RR^{i} \times \sC(Z),\RR^{i} \times \sC(Z)\big)\,\, \simeq \,\, \sO(i) \times \Aut\big(\sC(Z)\big).
    \end{equation}
\end{observation}

\begin{example}\label{ch5:: ex: codim 1 defect}
Since we are considering right fibrations with discrete fibers, via straightening, we can also interpret these as $\set$-valued presheaves on $\Bsc_{d}^{\cS_{d-1}}$. 

We would like to articulate the sense in which a $k$-dimensional submanifold inside a bigger $d$-dimensional manifold might carry an additional notion of \emph{color} or \emph{defect label} — this being a key feature in exploring how different submanifolds may themselves carry different theories. Let us define
\begin{equation}
    \cF\maps \big(\Bsc_{d}^{\cS_{d-1}}\big)^\op \longra \set
\end{equation}
by the rule
\begin{equation}
    \cF(U)=
    \begin{cases}
        A, & U \cong \RR^{k}\times \sC(S^{d-k-1}),\\
        *, & U \cong \RR^{d},\\
        \emptyset, & \text{otherwise,}
    \end{cases}
\end{equation}
where, as before, the link $S^{d-k-1}$ is purposely left unstratified. Define the associated discrete fibration as $\sD_{k \subset d}^{A} \longra \Bsc_{d}^{\cS_{d-1}}$.

Let $(M \ra P)$ be a stratified topological $d$-manifold. A $\sD_{k\subset d}^{A}$-framing amounts to a functor 
\begin{equation}
g \maps \entr(M) \longra \sD^{A}_{k\subset d}
\end{equation} 
such that the above diagram of $(\oo,1)$-categories commutes:
\begin{equation}
    \begin{tikzcd}[ampersand replacement=\&]
        \&\& {\sD^{A}_{k \subset d}} \\
        {\entr(M)} \&\& {\Bsc_d^{\cS_{d-1}}}\,\, .
        \arrow[from=1-3, to=2-3]
        \arrow["g", from=2-1, to=1-3]
        \arrow[from=2-1, to=2-3]
    \end{tikzcd}
\end{equation}
In particular, such a framing can exist only when every point of $M$ is locally modeled either on $\RR^d$ or on $\RR^{k}\times \sC(S^{d-k-1})$. Let $\gamma$ be a morphism in the enter-path category from the stratum $M_p$ to the stratum $M_q$, where the strata are labelled by connected components. If both $M_p$ and $M_q$ are smooth $d$-dimensional strata, then $g(\gamma)$ is determined; but if $M_q$ is a connected $k$-dimensional stratum, locally modeled on $\RR^{k}\times \sC(S^{d-k-1})$, then 
\begin{equation}
    g(\gamma) \maps * \longra A,
\end{equation}  
which is equivalent to a choice of $a\in A$ for every connected $k$-dimensional defect stratum. 
\end{example}



\begin{example}
    Let us now consider the following $2$-dimensional discrete defect datum.
    \begin{equation}
        \cF_{\DD} \maps (\Bsc_{2}^{\cS_1})^{\op} \longra \set,
    \end{equation}
    which, at the level of objects, maps:
    \begin{itemize}
        \item $\RR^{2}$ is assigned the set $D_2$.
        \item $\RR\times \sC(S^0)$ is assigned the set $D_1$.
        \item For each stratified topological $1$-sphere $\SS^1$, the basic $\sC(\SS^1)$ is assigned a set $D_0(\SS^1)$. This set should naturally carry the action of $\Aut(\SS^1)$; when $\SS^1$ is not the unstratified circle, this amounts to the cyclic permutations of the connected components. We write
        \begin{equation}
            D_0 \, \overset{\sf def}{=} \, \coprod_{[\SS^1]\in \pi_0(\cS_1)} D_0(\SS^1) 
        \end{equation}
    \end{itemize}
    At the level of morphisms, it acts as follows:
    \begin{itemize}
        \item There are two maps $\RR^2 \ra \RR\times \sC(S^0)$ which we will map to 
        \begin{equation}
            s,t \maps D_1 \longra D_2
        \end{equation}
        which we can equivalently conceive of as a single map 
        \begin{equation}
            (s,t) \maps D_1 \longra D_2 \times D_2.
        \end{equation} 
        \item For every isomorphism class of stratified topological 1-spheres $[\SS^1]\in \pi_0(\cS_1)$, there is a map 
        \begin{equation}
            D_0(\SS^1) \longra D_1,
        \end{equation} 
        for each $0$-dimensional stratum in $\SS^1$, and a map 
        \begin{equation}
            D_0(\SS^1) \longra D_2
        \end{equation}
        for each $1$-dimensional stratum in $\SS^1$. We depict these maps using the following commutative diagram induced by the entrance-paths of $\sC(\SS^1)$ with $\SS^1$ stratified by three $0$-dimensional strata.
        \begin{equation}    
            \tikzfig{figures/Stcob/SCO9}{\begin{tikzpicture}
	\begin{pgfonlayer}{nodelayer}
		\node [style=none] (0) at (-2.5, 0) {};
		\node [style=none] (1) at (2.5, 0) {};
		\node [style=none] (2) at (0, 0) {};
		\node [style=0 strat] (3) at (0, 0) {};
		\node [style=none] (4) at (-1.975, 1.7) {};
		\node [style=none] (5) at (0, -2.6) {};
		\node [style=none] (6) at (1.8, 1.825) {};
	\end{pgfonlayer}
	\begin{pgfonlayer}{edgelayer}
		\draw [style=defect patch 2] (1.center)
			 to [in=270, out=-90, looseness=1.75] (0.center)
			 to [in=90, out=90, looseness=1.75] cycle;
		\draw [style=ts-dashed, in=90, out=90, looseness=1.75] (0.center) to (1.center);
		\draw [style=ts-dashed, in=270, out=270, looseness=1.75] (1.center) to (0.center);
		\draw [style=geometric] (4.center) to (3);
		\draw [style=geometric] (3) to (6.center);
		\draw [style=geometric] (3) to (5.center);
	\end{pgfonlayer}
\end{tikzpicture}
} \qquad \xrightarrow{\,\, \mathsf{enter-paths}\,\,} \qquad \tikzfig{figures/Stcob/SCO17}{\begin{tikzpicture}
	\begin{pgfonlayer}{nodelayer}
		\node [style=none] (0) at (-2.5, 0) {};
		\node [style=none] (1) at (2.5, 0) {};
		\node [style=none] (2) at (0, 0) {};
		\node [style=0 strat] (3) at (0, 0) {};
		\node [style=none] (4) at (-1.975, 1.7) {};
		\node [style=none] (5) at (0, -2.6) {};
		\node [style=none] (6) at (1.8, 1.825) {};
		\node [style=dot] (7) at (-1.5, -1) {};
		\node [style=dot] (8) at (0, 2) {};
		\node [style=dot] (9) at (1.5, -1) {};
		\node [style=dot] (10) at (-1.5, 1.25) {};
		\node [style=dot] (11) at (1.475, 1.25) {};
		\node [style=dot] (12) at (0, -2) {};
		\node [style=none] (13) at (0, 2) {};
		\node [style=none] (14) at (0, 0.35) {};
		\node [style=none] (15) at (0, -2) {};
		\node [style=none] (16) at (0, 0) {};
		\node [style=none] (17) at (0, -0.25) {};
		\node [style=none] (18) at (0.25, -2) {};
		\node [style=none] (19) at (1.475, 1) {};
		\node [style=none] (20) at (-0.25, 2) {};
		\node [style=none] (21) at (-1.5, -0.75) {};
		\node [style=none] (22) at (0.25, 2) {};
		\node [style=none] (23) at (-0.25, -2) {};
		\node [style=none] (24) at (-0.275, -0.175) {};
		\node [style=none] (25) at (0.325, -0.15) {};
		\node [style=none] (26) at (-0.25, 0.25) {};
		\node [style=none] (27) at (0.25, 0.25) {};
		\node [style=col dot] (28) at (-1.5, 1.25) {};
		\node [style=col dot] (29) at (1.475, 1.25) {};
		\node [style=col dot] (30) at (0, -2) {};
		\node [style=none] (31) at (1.5, -0.7) {};
		\node [style=none] (32) at (1.25, -1.25) {};
		\node [style=none] (33) at (-1.25, -1.25) {};
	\end{pgfonlayer}
	\begin{pgfonlayer}{edgelayer}
		\draw [style=ts-dashed] (1.center)
			 to [in=270, out=-90, looseness=1.75] (0.center)
			 to [in=90, out=90, looseness=1.75] cycle;
		\draw [style=ts-dashed, in=90, out=90, looseness=1.75] (0.center) to (1.center);
		\draw [style=ts-dashed, in=270, out=270, looseness=1.75] (1.center) to (0.center);
		\draw [style=dir] (13.center) to (14.center);
		\draw [style=dir] (15.center) to (17.center);
		\draw [style=dir] (10) to (20.center);
		\draw [style=dir] (10) to (21.center);
		\draw [style=dir] (11) to (22.center);
		\draw [style=dir] (7) to (24.center);
		\draw [style=dir] (9) to (25.center);
		\draw [style=dir] (10) to (26.center);
		\draw [style=dir] (11) to (27.center);
		\draw [style=dir] (29) to (31.center);
		\draw [style=dir] (30) to (32.center);
		\draw [style=dir] (30) to (33.center);
	\end{pgfonlayer}
\end{tikzpicture}
} \qquad \xrightarrow{\,\, \cF_{\DD}\,\,} \qquad \tikzfig{figures/Stcob/SCO18}{\begin{tikzpicture}
	\begin{pgfonlayer}{nodelayer}
		\node [style=none] (2) at (0, 0) {};
		\node [style=0 strat] (3) at (0, 0) {};
		\node [style=none] (4) at (-1.975, 1.7) {$\scriptstyle{D_2}$};
		\node [style=none] (5) at (0, -2.6) {$\scriptstyle{D_2}$};
		\node [style=none] (6) at (1.8, 1.825) {$\scriptstyle{D_2}$};
		\node [style=dot] (7) at (-1.5, -1) {};
		\node [style=dot] (8) at (0, 2) {};
		\node [style=dot] (9) at (1.5, -1) {};
		\node [style=dot] (10) at (-1.5, 1.25) {};
		\node [style=dot] (11) at (1.475, 1.25) {};
		\node [style=dot] (12) at (0, -2) {};
		\node [style=none] (13) at (0, 2) {};
		\node [style=none] (14) at (0, 0.35) {};
		\node [style=none] (15) at (0, -2) {};
		\node [style=none] (16) at (0, 0) {};
		\node [style=none] (17) at (0, -0.25) {};
		\node [style=none] (20) at (-0.25, 2) {};
		\node [style=none] (21) at (-1.5, -0.75) {};
		\node [style=none] (22) at (0.25, 2) {};
		\node [style=none] (24) at (-0.275, -0.175) {};
		\node [style=none] (25) at (0.325, -0.15) {};
		\node [style=none] (26) at (-0.25, 0.25) {};
		\node [style=none] (27) at (0.25, 0.25) {};
		\node [style=col dot] (28) at (-1.5, 1.25) {};
		\node [style=col dot] (29) at (1.475, 1.25) {};
		\node [style=col dot] (30) at (0, -2) {};
		\node [style=none] (31) at (0, 2.5) {$\scriptstyle{D_1}$};
		\node [style=none] (32) at (-2, -1.25) {$\scriptstyle{D_1}$};
		\node [style=none] (33) at (2, -1.25) {$\scriptstyle{D_1}$};
		\node [style=none] (34) at (0.75, 0) {$\scriptscriptstyle{D_0}$};
		\node [style=none] (35) at (-0.95, 1.95) {$\scriptstyle{s}$};
		\node [style=none] (36) at (1.35, -1.2) {};
		\node [style=none] (37) at (-1.375, -1.2) {};
		\node [style=none] (38) at (1.5, -0.75) {};
		\node [style=none] (40) at (-1.05, -1.85) {$\scriptstyle{s}$};
		\node [style=none] (41) at (1.925, 0.2) {$\scriptstyle{s}$};
		\node [style=none] (42) at (-1.9, 0.2) {$\scriptstyle{t}$};
		\node [style=none] (43) at (1.05, -1.85) {$\scriptstyle{t}$};
		\node [style=none] (44) at (1.025, 1.975) {$\scriptstyle{t}$};
	\end{pgfonlayer}
	\begin{pgfonlayer}{edgelayer}
		\draw [style=dir] (13.center) to (14.center);
		\draw [style=dir] (15.center) to (17.center);
		\draw [style=dir] (10) to (20.center);
		\draw [style=dir] (10) to (21.center);
		\draw [style=dir] (11) to (22.center);
		\draw [style=dir] (7) to (24.center);
		\draw [style=dir] (9) to (25.center);
		\draw [style=dir] (10) to (26.center);
		\draw [style=dir] (11) to (27.center);
		\draw [style=dir] (30) to (36.center);
		\draw [style=dir] (30) to (37.center);
		\draw [style=dir] (29) to (38.center);
	\end{pgfonlayer}
\end{tikzpicture}
}
            \medskip
        \end{equation}
        These maps organize to a single map 
        \begin{equation}
            D_0 \longra \coprod_{[\SS^1]\in \pi_0(\cS_1)} \left(D_1 \times_{D_2} D_1 \times_{D_2} \cdots \times_{D_2} D_1 \right)
        \end{equation}
        akin to the ``fold map'' defined in \cite[Equation~2.27]{CRS19}. We justify the existence of the fold map by noting that it is given by the iterated product of the following dashed maps
        \begin{equation}
            \begin{tikzcd}[ampersand replacement=\&]
                {D_0} \\
                \& {D_1\times_{D_2}D_1} \&\& {D_1} \\
                \\
                \& {D_1} \&\& {D_2}\, .
                \arrow[dashed, from=1-1, to=2-2]
                \arrow[curve={height=-12pt}, from=1-1, to=2-4]
                \arrow[curve={height=12pt}, from=1-1, to=4-2]
                \arrow[from=2-2, to=2-4]
                \arrow[from=2-2, to=4-2]
                \arrow["s", from=2-4, to=4-4]
                \arrow["t"', from=4-2, to=4-4]
            \end{tikzcd}
        \end{equation}
    \end{itemize} 

    This presentation of the data involved in defining a discrete right fibration over $\Bsc_{2}^{\cS_{1}}$ is very much the unoriented analogue of the Carqueville--Runkel--Schaumann notion of $2$-dimensional defect data. In fact, one can say more: if each stratum of the basic was further endowed with an orientation, and we demanded all the structure functions to be equivariant with respect to this additional structure, we would recover their definition verbatim. It is with that goal in mind that we initiate our study of orientations for conically smooth stratified spaces.
\end{example}


\subsection{Orientations}

We end this section by sketching the expected correspondence between the $2$-dimensional defect data outlined in \cite[Definition~2.4]{CRS19} and $2$-dimensional discrete right fibrations over a suitably oriented version of $\Bsc_2^{\cS_1}$. The point is not to prove a theorem in a fully specified comparison category, but rather to explain why the right-fibration viewpoint appears to recover exactly the same local ingredients.

\begin{construction}
We are especially interested in studying the case where each stratum of our stratified $d$-manifold is equipped with an orientation. We demand no further compatibilities between these orientations. 

\medskip

\nid We construct these objects by using an appropriate notion of $\cB$-framing on the constructible tangent bundle of a stratified $d$-manifold. Let $\sB\sS\sO$ be the following space 
\begin{equation}
    \sB\sS\sO \, \overset{\sf def}{=} \, \bigcup_{d \geq 0} \sB\sS\sO(d)
\end{equation}
where $\sB\sS\sO(d)$ is the Grassmannian of $d$-planes in $\RR^{\oo}$ equipped with an orientation. The map which, for every $d\geq 0$, forgets the orientation 
\begin{equation}
    \sB\sS\sO(d) \longra \sB\sO(d)
\end{equation}
induces a map 
\begin{equation}
    \sB\sS\sO \longra \sB\sO.
\end{equation}
This is the stable tangential structure of orientations.
\medskip

\nid Recall that lifts
\begin{equation}
    \begin{tikzcd}[ampersand replacement=\&]
        \&\& {\sB\sS\sO} \\
        M \&\& {\sB\sO}
        \arrow[from=1-3, to=2-3]
        \arrow[from=2-1, to=1-3]
        \arrow["\tau"', from=2-1, to=2-3]
    \end{tikzcd}
\end{equation}
correspond to a choice of orientation for the vector bundle over $M$ which $\tau$ classifies. We can therefore package orientations into a tangential structure $\Vect^{\inj, \BSO}$, which we define as a modification of the original complete Segal space:

\smallskip

\begin{equation}
    [p] \longmapsto \Big| \Big\{V_0 \xlonghookra{f_1} V_1 \xlonghookra{f_2}\ldots \xlonghookra{f_p} V_p \text{ over } \AA^{\bullet} \Big\} \Big|
\end{equation}
\smallskip

\nid where each $V_{i} \ra \AA^{\bullet}$ is equipped with a map $\AA^{\bullet} \longra \BSO$ endowing it with an orientation. \emph{Crucially}, the only requirement is that fiberwise isomorphisms are orientation-preserving. The maximal $\oo$-subgroupoid of this $(\oo,1)$-category can be identified with  
\begin{equation}
    \bigsqcup_{k \geq 0} \BSO(k).
\end{equation}


\medskip

\nid A slightly more \emph{naive}, but useful, perspective on this $(\oo,1)$-category can be seen through its homotopy category. Its set of objects contains all finite-dimensional real vector spaces $V$, together with an $\sS\sO(k)$-torsor $X_k$ and an equivariant map $X_k \longra \mathsf{Fr}(V)$ to the space of frames, which amounts to a choice of orientation for $V$. Morphisms are injections of linear maps between them. That this is a complete Segal space follows arguments similar to those used to prove $\Vect^{\inj}$ is a complete Segal space.

\medskip

To bound the possible dimensions of the spaces we are working with, consider the pullback $(\oo,1)$-category
\begin{equation}
    \begin{tikzcd}[ampersand replacement=\&]
        {\Vect^{\inj,\BSO}_{\leq d}} \&\& {\Vect^{\inj,\BSO}} \\
        \\
        {\Vect^{\inj}_{\leq d}} \&\& {\Vect^{\inj}} \, .
        \arrow[from=1-1, to=1-3]
        \arrow[from=1-1, to=3-1]
        \arrow[from=1-3, to=3-3]
        \arrow[from=3-1, to=3-3]
    \end{tikzcd}
\end{equation}
which truncates the space of objects by cutting off the dimension of the vector bundles at $d$.

\nid To endow a conically smooth stratified space $M$ with a $\Vect^{\inj, \BSO}_{\leq d}$-framing is to endow every $k$-dimensional stratum $M_k$ with a lift $M_k \longra \sB\sS\sO(k)$ 

\begin{equation}
    \begin{tikzcd}[ampersand replacement=\&]
        \&\& {\sB\sS\sO(k)} \\
        M_{k} \&\& {\sB\sO(k)} \, .
        \arrow[from=1-3, to=2-3]
        \arrow[from=2-1, to=1-3]
        \arrow["\sT_{M_{k}}"', from=2-1, to=2-3]
    \end{tikzcd}
\end{equation}

\smallskip

\nid In particular, every connected stratum $M_k$ is endowed with an orientation and no further compatibility data. 

\nid Having stipulated the form of framing we are interested in, we want to consider stratified spaces equipped with such framings and constructible families thereof. To do that, we make use of the technology from Chapter 4. Let $\cB = \Vect^{\inj, \BSO}_{\leq d}$, we define the $(\oo,1)$-category of oriented stratified topological $d$-manifold basics as the following pullback $(\oo,1)$-category

\begin{equation}
    \begin{tikzcd}[ampersand replacement=\&]
        {\Bsc_{d}^{\cS_{d-1}, \mathsf{or}}} \&\& {\Bsc^{\cS_{d-1}}_{d}} \\
        \\
        {\Bun^{\cB,\emb}} \& {\Bun^{\emb}} \& {\Snglr = \Strat^{\emb}} 
        \arrow[from=1-1, to=1-3]
        \arrow[from=1-1, to=3-1]
        \arrow[hook, from=1-3, to=3-3]
        \arrow[from=3-1, to=3-2]
        \arrow["\simeq", from=3-2, to=3-3]
    \end{tikzcd}
\end{equation}
\end{construction}


\begin{observation}[A sketch of the oriented $2$-dimensional dictionary]
    Let \[
        \Bsc_{2}^{\cS_{1},\mathsf{or},\DD} \longra \Bsc_{2}^{\cS_{1}, \mathsf{or}}
    \]
    be a $2$-dimensional oriented discrete defect datum, and let
    \begin{equation}
        \cF \maps \big(\Bsc_{2}^{\cS_{1},\mathsf{or}}\big)^{\op} \longra \set
    \end{equation}
    be its straightening. Then one expects $\cF$ to determine, and to be determined by, defect data of the form
    \begin{equation}
        \DD = \big(D_0,D_1,D_2;f_0,f_1,f_2\big)
    \end{equation}
    in the sense of \cite{CRS19}.
\end{observation}

\medskip

\nid We will not attempt to formulate this as a theorem in a fully specified comparison category. Instead, we record the dictionary one reads off from the presheaf $\cF$.

\medskip

\nid First, the values of $\cF$ on the oriented local models supply the sets of labels:
\begin{enumerate}
    \item The value
    \begin{equation}
        D_2 \, \overset{\sf def}{=} \, \cF(\RR^{2})
    \end{equation}
    is the set of labels for oriented $2$-strata.
    \item The value
    \begin{equation}
        D_1 \, \overset{\sf def}{=} \, \cF\big(\RR\times \sC(S^0)\big)
    \end{equation}
    is the set of labels for oriented $1$-strata.
    \item For each oriented stratified $1$-sphere $\SS^1$, the value
    \begin{equation}
        D_0(\SS^1) \, \overset{\sf def}{=} \, \cF\big(\sC(\SS^1)\big)
    \end{equation}
    supplies the labels for oriented $0$-strata of singularity type $\SS^1$. Taking the coproduct over isomorphism classes gives
    \begin{equation}
        D_0 \, \overset{\sf def}{=} \, \coprod_{[\SS^1]} D_0(\SS^1),
    \end{equation}
    where the coproduct ranges over isomorphism classes of oriented stratified $1$-spheres.
\end{enumerate}

\nid Next, the morphisms in $\Bsc_{2}^{\cS_{1},\mathsf{or}}$ induce the familiar incidence maps. The two embeddings of an oriented codimension-$1$ basic into an oriented surface basic determine source and target maps
\begin{equation}
    s,t \maps D_1 \longra D_2,
\end{equation}
while, for each oriented stratified circle $\SS^1$, the entrance-paths in $\sC(\SS^1)$ induce maps from $D_0(\SS^1)$ to the adjacent copies of $D_1$ and $D_2$. These organize into the same kind of fold-type compatibility data that appears in \cite[Definition~2.4]{CRS19}.

\medskip

\nid Conversely, given $2$-dimensional oriented defect data in the sense of Carqueville--Runkel--Schaumann, one expects to reconstruct a set-valued presheaf on $\Bsc_{2}^{\cS_{1},\mathsf{or}}$ by assigning these sets to the corresponding oriented basics and using the structure maps to define restriction along morphisms. The point is that, once orientations are built into the local models, the ambiguity present in the unoriented codimension-$1$ case disappears: the two adjacent $2$-strata are no longer merely an unordered pair, but carry a source and a target.

\medskip

\nid In this way, the right-fibration viewpoint appears to recover precisely the same local combinatorics as the defect-data formalism of \cite{CRS19}. Making this into a genuine equivalence statement would require specifying an appropriate comparison category of oriented defect data and checking that the two constructions are inverse. We do not attempt that here.

\begin{observation}
We note that this discussion admits an immediate generalization to any other kind of stable tangential structure 
\begin{equation}
    \sB G \longra \BO.
\end{equation}
Likewise, the oriented $2$-dimensional dictionary sketched above should extend, at least formally, to higher dimensions: one expects the same right-fibration viewpoint to produce the relevant defect data by reading off labels from higher-dimensional oriented basics and the corresponding incidence maps from their entrance-path combinatorics. The obstruction is not so much conceptual as combinatorial. Already in dimension $2$, the bookkeeping becomes nontrivial; in higher dimensions, the resulting local configurations rapidly become harder to write down, visualize, and parse.

\medskip

\nid If we replace $\Vect^{\inj,\BSO}_{\le d}$ with $\Vect^{\inj,\sB G}_{\le d}$, then we obtain an $(\oo,1)$-category $\SMfd_{d}^{\sB G}$. If $G=*$, then $\SMfd_d^{\sB *}$ is the $(\oo,1)$-category of stratified topological $d$-manifolds whose $k$-strata admit a $k$-framing. 
\end{observation}

\begin{definition}
    Let $\sB G \longra \BO$ be a stable tangential structure. We define the $(\oo,1)$-category of $d$-dimensional $\sB G$-framings to be the pullback
    \begin{equation}
        \begin{tikzcd}[ampersand replacement=\&]
            {\Vect^{\inj,\sB G}_{\leq d}} \&\& {\Vect^{\inj, \sB G}} \\
            \\
            {\Vect^{\inj}_{\leq d}} \&\& {\Vect^{\inj}}.
            \arrow[from=1-1, to=1-3]
            \arrow[from=1-1, to=3-1]
            \arrow[from=1-3, to=3-3]
            \arrow[from=3-1, to=3-3]
        \end{tikzcd}
    \end{equation}
    Denoting this pullback $(\oo,1)$-category by $\cB$, we construct the $(\oo,1)$-category of stratified $d$-dimensional $\sB G$-basics to be the pullback $(\oo,1)$-category 
    \begin{equation}
        \begin{tikzcd}[ampersand replacement=\&]
            {\Bsc_{d}^{\cS_{d-1}, \sB G}} \&\& {\Bsc^{\cS_{d-1}}_{d}} \\
            \\
            {\Bun^{\cB,\emb}} \& {\Bun^{\emb}} \& {\Snglr = \Strat^{\emb}} \\
            {.}
            \arrow[from=1-1, to=1-3]
            \arrow[from=1-1, to=3-1]
            \arrow[hook, from=1-3, to=3-3]
            \arrow[from=3-1, to=3-2]
            \arrow["\simeq", from=3-2, to=3-3]
        \end{tikzcd}
    \end{equation}
\end{definition}


\section{Stratified topological quantum field theories}

\begin{definition}
    Let $\cC$ be a symmetric monoidal $(\oo,1)$-category. A \textbf{$d$-dimensional unoriented stratified topological quantum field theory} with values in $\cC$ is a symmetric monoidal functor 
    \begin{equation}
        \sZ \maps \bordstr_{d} \longra \cC.
    \end{equation}
\end{definition}

\begin{observation}
    There is an evident map of interval Segal spaces 
    \begin{equation}
        \bord_d \xhookrightarrow{\,\,\,\,\,\,\,\,} \bordstr_{d}
    \end{equation}
    which, for $I \in \Int$, sends a smooth submanifold 
    \begin{equation}
        \big( M\subset \AA^{\bullet}\times \core(I) \times \RR^{\oo}\big) 
    \end{equation}
    to the same manifold equipped with the trivial stratification. In that case, the constructible tangent bundle of $\big(M \ra *\big)$ is simply the ordinary tangent bundle of $M$, so the stratified tangential datum reduces to the usual rank-$d$ tangent bundle. In other words, one recovers no more data than that of a smooth $(d+\bullet)$-dimensional manifold with boundary.
\end{observation}

\begin{proposition}
    The above map induces a symmetric monoidal faithful functor
    \begin{equation}
        \ho(\bord_d) \xhookrightarrow{\,\,\,\,\,\,\,\,} \ho(\bordstr_{d}),
    \end{equation}
    which for $d>0$ is not full.
\end{proposition}
\begin{proof}
    Symmetric monoidality is immediate, since passing to the trivial stratification commutes with disjoint union.

    To prove faithfulness, we compare with the ordinary bordism categories. By \Cref{ch3:: thm: homotopy category bordism} and \Cref{ch5:: thm: homotopy category stratified bord}, the homotopy categories $\ho(\bord_d)$ and $\ho(\bordstr_d)$ are identified with $\mathsf{Cob}_d$ and $\cS\mathsf{Cob}_d$, respectively. Under these identifications, the induced functor sends the isomorphism class of a smooth bordism to the same bordism endowed with the trivial stratification. If two smooth bordisms become equal in $\cS\mathsf{Cob}_d$, then there is an isomorphism of stratified bordisms between their trivially stratified representatives; but such an isomorphism is exactly a diffeomorphism of smooth bordisms. Hence the original smooth bordism classes were already equal, so the functor is faithful.

    For $d>0$ the functor is not full. Indeed, the cylinder $S^{d-1}\times[0,1]$ admits nontrivial conically smooth stratifications with unstratified incoming and outgoing boundary, for example by declaring the interior hypersurface $S^{d-1}\times\{\frac12\}$ to be a lower-dimensional stratum. This determines a morphism in $\cS\mathsf{Cob}_d$ between unstratified closed $(d-1)$-manifolds which does not lie in the image of any smooth bordism class.
\end{proof}

\begin{observation}
    Every $d$-dimensional unoriented stratified topological quantum field theory 
    \begin{equation}
    \sZ \maps \bordstr_d \longra \cC
    \end{equation} with values in $\cC$ determines, by restriction, a $d$-dimensional unoriented topological quantum field theory with values in $\cC$ 
    \begin{equation}
        \sZ_{\sf sm} \maps \bord_d \xhookrightarrow{\,\,\,\,\,\,\,\,} \bordstr_d \xra{\,\,\,\, \sZ\,\,\,\,} \cC.
    \end{equation}
\end{observation}

More generally, one should expect many other full sub-$(\oo,1)$-categories obtained by imposing local singularity conditions encoded by a right fibration $\cB \ra \Bsc$, provided the allowed singularities retain spherical links of the appropriate codimension.

From this perspective, the present notion of bordism for conically smooth stratified spaces is flexible enough to accommodate examples beyond spaces with an underlying topological $n$-manifold. One might, for instance, try to incorporate finite directed graphs.

